% ch3.tex -- Module 6, Chapter 3: Electromagnetic induction (flux, Faraday, Lenz)
% Writer tasks: mod6-ch3-writer (first draft), mod6-ch3-finish (completion and layout),
%               mod6-ch3-editor (fixes after numeric, physics and novice reviews).
% Figure files: figures/ch3_falling.py (+ .pdf). All other figures are TikZ below.
\chapter{Electromagnetic Induction: Flux, Faraday and Lenz}\label{ch3:chap}
\covers{M6.IQ3.01, M6.IQ3.02}
% Local TikZ helpers (chapter prefix, see STYLE_GUIDE 1): a small open loop with
% one arrowhead that sits on the tangent, so the sense is clear even at small sizes.
% \chThreeACW{x}{y}{r} = anticlockwise loop, \chThreeCW{x}{y}{r} = clockwise loop.
\newcommand\chThreeACW[3]{%
  \draw[line width=1.2pt, hscink] ($(#1,#2)+(110:#3)$) arc (110:400:#3);
  \draw[line width=1.2pt, hscink, -{Latex[length=2.4mm, width=2.0mm]}]
    ($(#1,#2)+(40:#3)+(310:0.05)$) -- ($(#1,#2)+(40:#3)+(130:0.19)$);}
\newcommand\chThreeCW[3]{%
  \draw[line width=1.2pt, hscink] ($(#1,#2)+(70:#3)$) arc (70:-220:#3);
  \draw[line width=1.2pt, hscink, -{Latex[length=2.4mm, width=2.0mm]}]
    ($(#1,#2)+(140:#3)+(230:0.05)$) -- ($(#1,#2)+(140:#3)+(50:0.19)$);}

\inquiry{How are electric and magnetic fields related?}

\begin{goals}
  \item describe magnetic flux and the three ways it can change, using
        $\Phi = B_\parallel A = BA\cos\theta$
  \item calculate an induced emf and current with
        $\varepsilon = -N\dfrac{\Delta\Phi}{\Delta t}$ (Faraday's law)
  \item find the direction of an induced current (Lenz's law) and link it to
        conservation of energy
  \item explain induction in moving wires, solenoids, metal plates and two
        nearby coils
\end{goals}

\begin{syllabus}
  \item M6.IQ3.01: describe how magnetic flux can change, using
        $\Phi = B_\parallel A = BA\cos\theta$
  \item M6.IQ3.02: analyse, in words and with numbers, examples of Faraday's and
        Lenz's laws, $\varepsilon = -N\dfrac{\Delta\Phi}{\Delta t}$, including energy
        transfers: the emf, and the evidence for Lenz's law, when a magnet moves
        relative to straight wires, metal plates and solenoids; the emf in one
        solenoid caused by moving, or changing the current in, a nearby solenoid
\end{syllabus}

In 1820 Ørsted showed that a current makes a magnetic field. In 1831 Michael
Faraday showed the reverse: a magnet can make a current, but only while the
magnetism through the circuit is \emph{changing}. Generators, wireless chargers
and electric guitars all use this idea.

%% ===================================================================
\section{Magnetic flux}\label{ch3:sec:flux}
\covers{M6.IQ3.01}

\begin{bigidea}
Magnetic flux tells you how much magnetic field passes through a loop. It
depends on the field strength, the area and the angle. Change any one of them
and the flux changes.
\end{bigidea}

\begin{keywords}
  \kw{Magnetic flux ($\Phi$)}{how much magnetic field passes through an area.
    Unit: weber (Wb).}
  \kw{Magnetic field strength ($B$)}{how strong the field is. It is also called
    magnetic flux density. Unit: tesla (T).}
  \kw{Normal}{an imaginary arrow at \ang{90} to the flat surface of a loop.}
\end{keywords}

\begin{refresher}{Field lines, dots and crosses}
\begin{itemize}
  \item Field lines leave an N pole and enter an S pole. Closer lines mean a
        stronger field. A \term{uniform field} (same everywhere) has equally
        spaced, parallel lines.
  \item Dots (\,$\bullet$\,) show a field coming out of the page, like an arrow tip
        coming at you. Crosses ($\times$) show a field going into the page.
\end{itemize}
\end{refresher}

Hold a butterfly net in the wind. Facing the wind, lots of air goes through it.
Side-on, none does. \term{Magnetic flux} ($\Phi$, capital ``phi'') is the same idea
for a magnetic field: it is how much field passes through an area. Picture it as
the number of field lines that pass through the loop.

To describe how a loop faces the field, draw its \term{normal}. This is an arrow at
\ang{90} to the surface. The angle $\theta$ is measured between $B$ and the normal,
\emph{not} between $B$ and the surface (\cref{ch3:fig:angle}).

The \term{plane} of a loop is the flat surface inside it, like the skin of a drum.
If the field is \emph{perpendicular} (at \ang{90}) to the plane, it goes straight
through: $\theta = \ang{0}$. If the field is \emph{parallel} to the plane, it
slides along the surface: $\theta = \ang{90}$.

\begin{toolkit}{Areas, units and $\cos\theta$}
\begin{itemize}
  \item Splitting $B$ (like splitting a vector in \cref{ch2:chap}): when $\theta$ is
        measured from the normal, the part of $B$ along the normal is
        $B_\parallel = B\cos\theta$. At $\theta = \ang{0}$ all of $B$ goes through
        ($\cos\ang{0} = 1$). At $\theta = \ang{90}$ none goes through
        ($\cos\ang{90} = 0$). In $F = lIB\sin\theta$ the angle is measured
        differently, so do not mix them up.
  \item Rectangle: $A = \text{length} \times \text{width}$. Circle: $A = \pi r^2$, with
        $r$ the radius (half the diameter).
  \item Change lengths to metres first. $\qty{1}{cm^2} = \qty{e-4}{m^2}$, so
        $\qty{12}{cm} \times \qty{8.0}{cm} = \qty{0.12}{m} \times \qty{0.080}{m} = \qty{9.6e-3}{m^2}$.
  \item $\qty{1}{mT} = \qty{e-3}{T}$, \ $\qty{1}{\micro T} = \qty{e-6}{T}$, \
        $\qty{1}{mWb} = \qty{e-3}{Wb}$.
  \item $\cos\ang{0} = 1$, \ $\cos\ang{90} = 0$ (calculator in degree mode). Angle
        given to the \emph{surface}? Use $\theta = \ang{90} - \text{that angle}$.
\end{itemize}
\end{toolkit}

\begin{figure}
  \centering
  \begin{tikzpicture}[x=1cm, y=1cm]
    % ---------- (a) theta = 0
    \begin{scope}
      \foreach \y in {0.2, 0.75, 1.8, 2.35} \draw[field] (0,\y) -- (3.6,\y);
      \node[diagramlabel, text=hscfield, right] at (3.6,2.35) {$B$};
      \coordinate (C) at (1.5,1.275);
      \draw[line width=2.4pt, hscbrand] ($(C)+(90:0.7)$) -- ($(C)+(270:0.7)$);
      \draw[->, line width=1.1pt, hscink] (C) -- ++(0:1.0)
        node[pos=1, right, font=\footnotesize, inner sep=2pt] {normal};
      \node[align=center, font=\small] at (1.8,-0.5)
        {(a) $\theta = \ang{0}$\\ $\Phi = BA$ (largest)};
    \end{scope}
    % ---------- (b) theta = 50 deg
    \begin{scope}[xshift=4.75cm]
      \foreach \y in {0.2, 0.75, 1.8, 2.35} \draw[field] (0,\y) -- (3.6,\y);
      \node[diagramlabel, text=hscfield, right] at (3.6,2.35) {$B$};
      \coordinate (C) at (1.5,1.275);
      \draw[guide] (C) -- ++(0:1.1);
      \draw[line width=2.4pt, hscbrand] ($(C)+(140:0.7)$) -- ($(C)+(320:0.7)$);
      \draw[->, line width=1.1pt, hscink] (C) -- ++(50:0.9)
        node[anchor=south west, font=\footnotesize, inner sep=1pt] {normal};
      \draw[line width=0.7pt, hscink] ($(C)+(0:0.45)$) arc (0:50:0.45);
      \node[font=\small, inner sep=0.5pt] at ($(C)+(22:0.68)$) {$\theta$};
      \draw[line width=0.8pt, hscmuted] ($(C)+(320:0.55)$) arc (320:360:0.55);
      \node[font=\footnotesize, text=hscmuted, inner sep=1pt, anchor=west]
        at ($(C)+(338:0.8)$) {$\ang{90} - \theta$};
      \node[align=center, font=\small] at (1.8,-0.5)
        {(b) tilted\\ $\Phi = BA\cos\theta$};
    \end{scope}
    % ---------- (c) theta = 90 deg
    \begin{scope}[xshift=9.5cm]
      \foreach \y in {0.2, 0.75, 1.8, 2.35} \draw[field] (0,\y) -- (3.6,\y);
      \node[diagramlabel, text=hscfield, right] at (3.6,2.35) {$B$};
      \coordinate (C) at (1.2,1.275);
      \draw[line width=2.4pt, hscbrand] ($(C)+(180:0.7)$) -- ($(C)+(0:0.7)$);
      \draw[->, line width=1.1pt, hscink] (C) -- ++(90:0.9)
        node[pos=0.8, anchor=east, font=\footnotesize, inner sep=2pt] {normal};
      \node[align=center, font=\small] at (1.8,-0.5)
        {(c) $\theta = \ang{90}$: plane\\ parallel to $B$, \ $\Phi = 0$};
    \end{scope}
  \end{tikzpicture}
  \caption{A flat loop in a uniform field, seen from the side (thick navy
    line).}\label{ch3:fig:angle}
\end{figure}

Only the part of $B$ along the normal, $B_\parallel = B\cos\theta$, passes through
the loop. When the loop faces the field square-on ($\theta = \ang{0}$), the flux
is largest. When its plane is parallel to $B$ ($\theta = \ang{90}$), the flux is
zero.

Flux and field strength are linked by their units: $\qty{1}{Wb} = \qty{1}{T.m^2}$.
When $B$ is perpendicular to the plane of the loop ($\theta = \ang{0}$),
$B = \dfrac{\Phi}{A}$: the flux per square metre. That is why $B$ is also called
the \term{magnetic flux density}.

\begin{formulacard}{Magnetic flux}
  \cardformula{\Phi = B_\parallel A = BA\cos\theta}
  \begin{symbols}
    \sym{$\Phi$}{magnetic flux through the loop}{\si{Wb} ($= \si{T.m^2}$)}
    \sym{$B_\parallel$}{the part of $B$ along the normal of the loop}{\si{T}}
    \sym{$B$}{magnetic field strength (flux density)}{\si{T}}
    \sym{$A$}{area of the loop that is inside the field}{\si{m^2}}
    \sym{$\theta$}{angle between $B$ and the \emph{normal} to the loop}{\si{\degree}
      (SI unit: rad)}
  \end{symbols}
  \inwords{flux = field strength $\times$ area $\times$ how squarely the loop
    faces the field.}
  \whentouse{a flat loop or coil in a uniform field.}
  \whennottouse{do not put in the angle to the surface (the plane). Change it to
    the angle to the normal first. $\Phi$ is the flux through \emph{one} turn. The
    number of turns $N$ goes into Faraday's law, not here.}
  \rearranged{$B = \dfrac{\Phi}{A\cos\theta}$ \qquad $A = \dfrac{\Phi}{B\cos\theta}$}
  \tensecond{$B = \qty{0.50}{T}$, $A = \qty{0.020}{m^2}$, $\theta = \ang{0}$: \
    $\Phi = 0.50 \times 0.020 \times 1 = \qty{1.0e-2}{Wb}$.}
  \trap{a loop whose plane is \emph{parallel} to the field has $\theta = \ang{90}$,
    so $\Phi = 0$. Change \si{cm^2} to \si{m^2} ($\times\,10^{-4}$).}
\end{formulacard}

\begin{workedexample}{Flux through a tilted loop}\label{ch3:we:tilted}
A square loop with \qty{15}{cm} sides is in a uniform \qty{0.32}{T} field.
Find the flux when (a) the field is perpendicular to the plane of the loop, and
(b) the field makes \ang{50} with the plane of the loop.
\given{side $= \qty{15}{cm} = \qty{0.15}{m}$, so $A = (\qty{0.15}{m})^2
  = \qty{0.0225}{m^2}$, \ $B = \qty{0.32}{T}$.\\
  (a) $B$ is along the normal: $\theta = \ang{0}$. \
  (b) \ang{50} to the plane: $\theta = \ang{90} - \ang{50} = \ang{40}$.}
\find{the flux $\Phi$ in each case}
\formula{$\Phi = BA\cos\theta$}
\substitute{\eqlines{\text{(a)}\ \Phi &= \qty{0.32}{T} \times \qty{0.0225}{m^2}
    \times \cos\ang{0} = \qty{7.2e-3}{Wb}\\
  \text{(b)}\ \Phi &= \qty{0.32}{T} \times \qty{0.0225}{m^2} \times \cos\ang{40}\\
    &= \qty{7.2e-3}{Wb} \times 0.766 = \qty{5.52e-3}{Wb}}}
\answer{(a) \final{$\qty{7.2e-3}{Wb}$} \quad (b) \final{$\qty{5.5e-3}{Wb}$}
  \ (2 significant figures)}
\checkstep{the tilted loop has less flux: fewer lines pass through it \cmark.
  Using $\cos\ang{50}$ by mistake would give \qty{4.6e-3}{Wb}.}
\end{workedexample}

\subsection{Three ways to change the flux}

Because $\Phi = BA\cos\theta$, the flux changes when $B$, $A$ or $\theta$ changes
(\cref{ch3:fig:ways}).
\begin{enumerate}
  \item \textbf{Change $B$}: move a magnet closer or further away; change the
        current in an electromagnet (a coil that acts as a magnet while a current
        flows); switch a field on, off, or reverse it.
  \item \textbf{Change $A$}: slide the loop into or out of the field; squash or
        stretch it; slide a rod along rails (\cref{ch3:sec:motional}).
  \item \textbf{Change $\theta$}: rotate the loop, as in a generator (\cref{ch5:chap}).
\end{enumerate}

\begin{figure}
  \centering
  \begin{tikzpicture}[x=1cm, y=1cm]
    % ---------- (a) change B
    \begin{scope}
      \fill[hscvel!22] (0,0.9) rectangle (0.7,1.4);
      \fill[hscforce!22] (0.7,0.9) rectangle (1.4,1.4);
      \draw[line width=0.8pt, hscink] (0,0.9) rectangle (1.4,1.4);
      \node[font=\small\bfseries] at (0.35,1.15) {S};
      \node[font=\small\bfseries] at (1.05,1.15) {N};
      \draw[velocity] (0.2,1.75) -- (1.2,1.75) node[right, text=hscvel] {$v$};
      \draw[field] (1.4,1.15) -- (3.6,1.15);
      \draw[field] (1.4,1.32) .. controls (2.2,1.75) .. (3.6,2.15);
      \draw[field] (1.4,0.98) .. controls (2.2,0.55) .. (3.6,0.15);
      \draw[line width=1.6pt, hscbrand] (2.75,1.15) ellipse (0.13 and 0.6);
      \node[align=center, text width=4.3cm, font=\small] at (1.8,-0.4)
        {(a) Change $B$: the magnet comes closer.};
    \end{scope}
    % ---------- (b) change A
    \begin{scope}[xshift=4.8cm]
      \draw[hscfield, dashed, line width=0.6pt] (0,0.2) rectangle (2.2,2.1);
      \foreach \x in {0.3, 0.75, 1.2} \foreach \y in {0.45, 0.9, 1.35, 1.8}
        \node[text=hscfield, font=\small] at (\x,\y) {$\times$};
      \foreach \x in {1.65, 2.05} \foreach \y in {0.45, 0.9, 1.35, 1.8}
        \node[text=hscfield, font=\small] at (\x,\y) {$\times$};
      \draw[line width=1.6pt, hscbrand] (1.45,0.65) rectangle (2.85,1.6);
      \draw[velocity] (2.85,1.125) -- (3.75,1.125) node[above, text=hscvel] {$v$};
      \node[align=center, text width=4.3cm, font=\small] at (1.8,-0.4)
        {(b) Change $A$: less of the loop is in the field.};
    \end{scope}
    % ---------- (c) change theta
    \begin{scope}[xshift=9.6cm]
      \foreach \y in {0.2, 0.85, 1.5, 2.15} \draw[field] (0,\y) -- (3.6,\y);
      \coordinate (C) at (1.8,1.175);
      \draw[line width=2.4pt, hscbrand] ($(C)+(115:0.8)$) -- ($(C)+(295:0.8)$);
      \fill[hscink] (C) circle (1.6pt);
      \draw[->, line width=0.9pt, hscmuted] ($(C)+(70:1.0)$) arc (70:20:1.0);
      \draw[->, line width=0.9pt, hscmuted] ($(C)+(250:1.0)$) arc (250:200:1.0);
      \node[align=center, text width=4.3cm, font=\small] at (1.8,-0.4)
        {(c) Change $\theta$: the loop turns.};
    \end{scope}
  \end{tikzpicture}
  \caption{Three ways to change the flux through a loop.}\label{ch3:fig:ways}
\end{figure}

A loop that slides sideways \emph{inside} a large uniform field keeps the same $B$,
$A$ and $\theta$. Its flux does not change, so no emf (voltage) is made
in it (\cref{ch3:sec:faraday}). Moving is not enough.

Flux can be positive or negative. The normal arrow points out of one face of the
loop. If the field goes through the loop the same way as the normal, the flux is
positive. If it goes through the opposite way, the flux is negative. Suppose the
field starts along the normal ($\theta = \ang{0}$, so $\Phi = +BA$). Reversing
the field (or turning the loop over) makes $\theta = \ang{180}$, and
$\cos\ang{180} = -1$. So the flux goes from $+BA$ to $-BA$. The size of the change
is $2BA$, not zero.

\begin{workedexample}{Changing the flux}\label{ch3:we:changeflux}
A \qty{20}{cm} by \qty{10}{cm} loop is in a \qty{0.60}{T} field, with its normal
at \ang{30} to the field. Find the starting flux. Then find the change in flux
when (a) the loop is turned until its plane is parallel to the field, or
(b) the field is reversed instead.
\given{$A = \qty{0.20}{m} \times \qty{0.10}{m} = \qty{0.020}{m^2}$, \
  $B = \qty{0.60}{T}$, \ $\theta_\mathrm{i} = \ang{30}$;\\
  (a) plane parallel to the field: $\theta_\mathrm{f} = \ang{90}$; \
  (b) field reversed: the normal is now at
  $\theta_\mathrm{f} = \ang{180} - \ang{30} = \ang{150}$ to the field
  ($\cos\ang{150} = -0.866$)}
\find{$\Phi_\mathrm{i}$, then $\Delta\Phi = \Phi_\mathrm{f} - \Phi_\mathrm{i}$ for (a) and (b)}
\formula{$\Phi = BA\cos\theta$, \quad $\Delta\Phi = \Phi_\mathrm{final} - \Phi_\mathrm{initial}$}
\substitute{\eqlines{\Phi_\mathrm{i} &= \qty{0.60}{T} \times \qty{0.020}{m^2} \times
    \cos\ang{30} = \qty{1.039e-2}{Wb}\\
  \text{(a)}\ \Phi_\mathrm{f} &= \qty{0.60}{T} \times \qty{0.020}{m^2} \times
    \cos\ang{90} = 0\\
  \Delta\Phi &= 0 - \qty{1.039e-2}{Wb} = \qty{-1.04e-2}{Wb}\\
  \text{(b)}\ \Phi_\mathrm{f} &= \qty{0.60}{T} \times \qty{0.020}{m^2} \times
    \cos\ang{150} = \qty{-1.039e-2}{Wb}\\
  \Delta\Phi &= (\qty{-1.039e-2}{Wb}) - (\qty{1.039e-2}{Wb})\\
  &= \qty{-2.08e-2}{Wb}}}
\answer{\final{$\Phi_\mathrm{i} = \qty{1.0e-2}{Wb}$}; \ (a) \final{$\qty{-1.0e-2}{Wb}$};
  \ (b) \final{$\qty{-2.1e-2}{Wb}$}}
\checkstep{(b) is twice as big as (a): $2 \times 1.04 = 2.08$ \cmark. In (a) the
  flux only falls to zero. In (b) it falls through zero to
  $-\qty{1.0e-2}{Wb}$.}
\end{workedexample}

\begin{mistakes}
\begin{itemize}
  \item \xmark\ Putting the angle to the surface into $\cos\theta$.
        \cmark\ $\theta = \ang{90} - \text{angle to the surface}$.
  \item \xmark\ Using the diameter in $\pi r^2$, or leaving $A$ in \si{cm^2}.
        \cmark\ Halve the diameter; multiply \si{cm^2} by $10^{-4}$.
  \item \xmark\ ``Reversing the field gives $\Delta\Phi = 0$.''
        \cmark\ $\lvert\Delta\Phi\rvert = 2BA$.
  \item \xmark\ Mixing up flux (\si{Wb}) and flux density (\si{T}).
        \cmark\ Flux $\Phi$ is the total field through the loop (\si{Wb}). Flux
        density $B$ is the field strength: the flux per square metre (\si{T}).
\end{itemize}
\end{mistakes}

\begin{examtip}
``Describe how the flux through a coil can change'' (2--3 marks): write
$\Phi = BA\cos\theta$, then give one example each for changing $B$, $A$ and
$\theta$. ``Contrast flux and flux density'' (2 marks): one mark for each meaning
with its unit, joined by a contrast word such as ``whereas''.
\end{examtip}

\begin{quickcheck}
  \item A circular loop of radius \qty{3.0}{cm} is in a \qty{0.20}{T} field that is
        perpendicular to the plane of the loop. Calculate the flux through
        it.\label{ch3:qc:circle}
  \item The loop in \cref{ch3:qc:circle} is turned so that the field makes \ang{25}
        with the plane of the loop. Calculate the new flux.
  \item Contrast magnetic flux and magnetic flux density. (2 marks)
  \item Loop X has area \qty{0.050}{m^2} and faces a \qty{0.20}{T} field
        square-on. Loop Y has area \qty{0.010}{m^2} and faces a \qty{0.50}{T}
        field square-on. (a) Calculate the flux through each loop. (b) A student
        says: ``More flux means a stronger field.'' Use your answers to evaluate
        this claim.
\end{quickcheck}

%% ===================================================================
\section{Faraday's law}\label{ch3:sec:faraday}
\covers{M6.IQ3.02}

\begin{bigidea}
When the flux through a coil changes, an emf is induced in it. The faster the
flux changes, and the more turns the coil has, the bigger the emf. A flux that
does not change induces nothing, however big it is.
\end{bigidea}

\begin{keywords}
  \kw{Induce}{to make something happen without touching it. An \emph{induced}
    emf is an emf made by a changing flux.}
  \kw{Electromagnetic induction}{making an emf by changing the magnetic flux
    through a circuit.}
  \kw{emf ($\varepsilon$)}{short for electromotive force. It is not a force: it
    is the energy a source gives to each coulomb of charge. Unit: volt (V).}
  \kw{Induced current}{the current that an induced emf drives around a closed
    circuit.}
\end{keywords}

Join a coil to a \term{galvanometer}, a sensitive meter whose needle can move
either way from zero (\cref{ch3:fig:galvo}). Push a magnet into the coil: the
needle flicks. Hold the magnet still: the needle reads zero. Pull the magnet out:
the needle flicks the other way. Move it faster, or use a stronger magnet or more
turns: the flicks get bigger.

\begin{figure}
  \centering
  \begin{tikzpicture}[x=1cm, y=1cm]
    \foreach \xs/\nang/\lab in {0/125/{(a) Push in: the needle\\ flicks left.},
        4.9/90/{(b) Hold still: the\\ needle reads zero.},
        9.8/55/{(c) Pull out: the needle\\ flicks right.}} {
    \begin{scope}[xshift=\xs cm]
      % magnet (S on the left, N on the right, facing the coil)
      \fill[hscvel!22] (0,0) rectangle (0.6,0.45);
      \fill[hscforce!22] (0.6,0) rectangle (1.2,0.45);
      \draw[line width=0.8pt, hscink] (0,0) rectangle (1.2,0.45);
      \node[font=\footnotesize\bfseries] at (0.3,0.225) {S};
      \node[font=\footnotesize\bfseries] at (0.9,0.225) {N};
      % coil
      \foreach \x in {1.75, 1.97, ..., 2.7}
        \draw[line width=1pt, hscink] (\x,0.225) ellipse (0.09 and 0.5);
      % leads and galvanometer
      \draw[line width=0.9pt, hscink] (1.75,-0.28) -- (1.75,-0.95) -- (1.93,-0.95);
      \draw[line width=0.9pt, hscink] (2.63,-0.28) -- (2.63,-0.95) -- (2.53,-0.95);
      \draw[line width=0.9pt, hscink, fill=white] (2.23,-0.95) circle (0.3);
      \draw[line width=0.5pt, hscmuted] ($(2.23,-1.08)+(150:0.3)$) arc (150:30:0.3);
      \draw[line width=1.1pt, hscforce] (2.23,-1.08) -- ++(\nang:0.33);
      \fill[hscink] (2.23,-1.08) circle (1.2pt);
      \node[font=\scriptsize, anchor=north, inner sep=1pt] at (2.23,-1.27) {G};
      \node[font=\footnotesize, align=center, anchor=north] at (1.6,-1.6) {\lab};
    \end{scope}}
    \draw[velocity] (0.1,0.8) -- (1.1,0.8) node[right, text=hscvel] {$v$};
    \node[font=\footnotesize, text=hscvel] at (5.5,0.8) {$v = 0$};
    \draw[velocity] (10.9,0.8) -- (9.9,0.8) node[left, text=hscvel] {$v$};
  \end{tikzpicture}
  \caption{A magnet and a coil joined to a centre-zero galvanometer (G). A
    current flows only while the magnet moves.}\label{ch3:fig:galvo}
\end{figure}

So an emf appears only while the flux is \emph{changing}. We say the changing
flux \term{induces} the emf: it causes the emf without any battery or contact.
Making an emf this way is called \term{electromagnetic induction}.

Each turn of the coil gets the same emf. The turns are joined one after another,
like cells in a battery, so their emfs add up. A coil of $N$ turns gives $N$ times
the emf of one turn. A wireless phone charger uses this idea: a changing current
in the pad changes the flux through a coil inside the phone.

\begin{refresher}{emf, potential difference and $V = IR$}
\begin{itemize}
  \item A source (battery, or a coil with changing flux) gives each coulomb of
        charge energy: its emf $\varepsilon$. A lamp or resistor takes energy:
        the potential difference $V$ across it.
  \item In a closed circuit, $I = \dfrac{\varepsilon}{R}$, with $R$ the
        \emph{total} resistance. With a gap in the circuit, the emf is still
        there, but no current flows.
\end{itemize}
\end{refresher}

\begin{formulacard}{Faraday's law of induction}
  \cardformula{\varepsilon = -N\frac{\Delta\Phi}{\Delta t}}
  \begin{symbols}
    \sym{$\varepsilon$}{average induced emf}{\si{V}}
    \sym{$N$}{number of turns in the coil}{no unit}
    \sym{$\Delta\Phi$}{change in flux through one turn ($\Delta$, ``delta'', means
      ``change in''): $\Phi_\mathrm{final} - \Phi_\mathrm{initial}$}{\si{Wb}}
    \sym{$\Delta t$}{time taken for the change}{\si{s}}
  \end{symbols}
  \inwords{emf = number of turns $\times$ how fast the flux through each turn
    changes. The minus sign means the emf opposes the change (Lenz's law,
    \cref{ch3:sec:lenz}).}
  \whentouse{the flux through a coil changes for any reason ($B$, $A$ or $\theta$).
    For a closed circuit, combine with $V = IR$ (on the sheet):
    $I = \dfrac{\varepsilon}{R}$, with $R$ the total resistance.}
  \whennottouse{to find which way the current flows round the coil. Use this
    formula for the size of the emf, and Lenz's law for the direction. (On an
    emf--time graph the sign only shows that two emfs are opposite.) It gives the
    \emph{average} emf. For the emf at one instant, use the gradient of the
    $\Phi$--$t$ graph.}
  \rearranged{$\Delta t = \dfrac{N\lvert\Delta\Phi\rvert}{\lvert\varepsilon\rvert}$
    \qquad $\lvert\Delta\Phi\rvert = \dfrac{\lvert\varepsilon\rvert\,\Delta t}{N}$
    \qquad $N = \dfrac{\lvert\varepsilon\rvert\,\Delta t}{\lvert\Delta\Phi\rvert}$}
  \tensecond{$N = 100$, $\Delta\Phi = \qty{3.0e-3}{Wb}$ in \qty{0.10}{s}: \
    $\lvert\varepsilon\rvert = \dfrac{100 \times \num{3.0e-3}}{0.10} = \qty{3.0}{V}$.}
  \trap{$\Delta\Phi$ is the \emph{change} in flux. A huge but steady flux gives
    zero emf. Use $N$ once only. Units: $\si{Wb.s^{-1}} = \si{V}$.}
\end{formulacard}

\begin{workedexample}{A field that grows}\label{ch3:we:growB}
A 250-turn coil, each turn of area \qty{12}{cm^2}, has its plane perpendicular to
a field ($\theta = \ang{0}$).
The field rises steadily from \qty{0.050}{T} to \qty{0.35}{T} in \qty{0.15}{s}.
(a) Find the size of the average emf. (b) The circuit's total resistance is
\qty{5.0}{\ohm}. Find the induced current.
\given{$N = 250$, \ $A = \qty{12e-4}{m^2} = \qty{1.2e-3}{m^2}$, \ $\theta = \ang{0}$,
  \ $B$: \qty{0.050}{T} $\to$ \qty{0.35}{T}, \ $\Delta t = \qty{0.15}{s}$, \
  $R = \qty{5.0}{\ohm}$}
\find{the size of $\varepsilon$, then $I$}
\formula{$\Phi = BA\cos\theta$, \quad
  $\lvert\varepsilon\rvert = N\dfrac{\lvert\Delta\Phi\rvert}{\Delta t}$, \quad
  $I = \dfrac{\varepsilon}{R}$}
\substitute{\eqlines{\Phi_\mathrm{i} &= \qty{0.050}{T} \times \qty{1.2e-3}{m^2}
    = \qty{6.0e-5}{Wb} && \text{start}\\
  \Phi_\mathrm{f} &= \qty{0.35}{T} \times \qty{1.2e-3}{m^2} = \qty{4.2e-4}{Wb}
    && \text{end}\\
  \Delta\Phi &= \qty{4.2e-4}{Wb} - \qty{6.0e-5}{Wb} = \qty{3.6e-4}{Wb}\\
  \lvert\varepsilon\rvert &= \frac{250 \times \qty{3.6e-4}{Wb}}{\qty{0.15}{s}}
    = \qty{0.60}{V}\\
  I &= \frac{\qty{0.60}{V}}{\qty{5.0}{\ohm}} = \qty{0.12}{A}}}
\answer{(a) \final{$\varepsilon = \qty{0.60}{V}$} \quad (b) \final{$I = \qty{0.12}{A}$}}
\checkstep{$\si{Wb.s^{-1}} = \si{V}$ \cmark. A fraction of a volt is typical for a
  school coil.}
\end{workedexample}

\begin{workedexample}{Turning a coil}\label{ch3:we:rotate}
An 80-turn coil, \qty{5.0}{cm} by \qty{4.0}{cm}, has its plane perpendicular to a
\qty{0.25}{T} field. (a) It is turned through \ang{90} in \qty{0.020}{s}. Find the
average emf. (b) A student says: ``Turning it through \ang{180} in \qty{0.040}{s}
doubles the emf, because the flux change doubles.'' Evaluate this claim.
\given{$N = 80$, \ $A = \qty{0.050}{m} \times \qty{0.040}{m} = \qty{2.0e-3}{m^2}$, \
  $B = \qty{0.25}{T}$, \ $\theta_\mathrm{i} = \ang{0}$; \ (a) $\theta_\mathrm{f} = \ang{90}$; \
  (b) $\theta_\mathrm{f} = \ang{180}$}
\find{the size of the average emf in (a) and (b)}
\formula{$\Phi = BA\cos\theta$, \quad
  $\lvert\varepsilon\rvert = N\dfrac{\lvert\Delta\Phi\rvert}{\Delta t}$}
\substitute{\eqlines{\Phi_\mathrm{i} &= \qty{0.25}{T} \times \qty{2.0e-3}{m^2} \times
    \cos\ang{0} = \qty{5.0e-4}{Wb}\\
  \text{(a)}\ \Phi_\mathrm{f} &= \qty{0.25}{T} \times \qty{2.0e-3}{m^2} \times
    \cos\ang{90} = 0\\
  \lvert\Delta\Phi\rvert &= \lvert 0 - \qty{5.0e-4}{Wb} \rvert
    = \qty{5.0e-4}{Wb}\\
  \lvert\varepsilon\rvert &= \frac{80 \times \qty{5.0e-4}{Wb}}{\qty{0.020}{s}}
    = \qty{2.0}{V}\\
  \text{(b)}\ \Phi_\mathrm{f} &= \qty{0.25}{T} \times \qty{2.0e-3}{m^2} \times
    \cos\ang{180} = \qty{-5.0e-4}{Wb}\\
  \lvert\Delta\Phi\rvert &= \lvert \qty{-5.0e-4}{Wb} - \qty{5.0e-4}{Wb} \rvert
    = \qty{1.0e-3}{Wb}\\
  \lvert\varepsilon\rvert &= \frac{80 \times \qty{1.0e-3}{Wb}}{\qty{0.040}{s}}
    = \qty{2.0}{V}}}
\answer{(a) \final{$\varepsilon = \qty{2.0}{V}$} \quad (b) also \qty{2.0}{V}, so the
  claim is \textbf{incorrect}: the flux change doubles, but so does the time.}
\checkstep{the rate of change of flux is the same in both:
  $\qty{5.0e-4}{Wb} \div \qty{0.020}{s} = \qty{1.0e-3}{Wb} \div \qty{0.040}{s}
  = \qty{0.025}{Wb.s^{-1}}$, so the emf is the same \cmark. These are
  \emph{average} values; the emf varies while the coil turns.}
\end{workedexample}

\begin{toolkit}{The gradient of a graph}
\begin{itemize}
  \item Gradient $= \dfrac{\text{rise}}{\text{run}} = \dfrac{\Delta y}{\Delta x}$,
        from two points on a straight section.
  \item On a $\Phi$--$t$ graph the gradient is $\dfrac{\Delta\Phi}{\Delta t}$
        (\si{Wb.s^{-1}}), so $\lvert\varepsilon\rvert = N \times \text{gradient}$.
  \item Steep: big emf. Flat: zero emf. Rising and falling parts give opposite
        signs.
  \item For a curve, draw a \term{tangent} (a straight line that just touches the
        curve at one point) and find its gradient.
  \item Change prefixes first: \qty{1}{mWb} $= \qty{e-3}{Wb}$, \
        \qty{1}{ms} $= \qty{e-3}{s}$.
\end{itemize}
\end{toolkit}

\begin{workedexample}{Reading a flux--time graph}\label{ch3:we:graph}
\Cref{ch3:fig:phigraph}(a) shows the flux through each turn of a 50-turn coil.
Find the emf in each part, and draw the emf--time graph.
\given{$N = 50$; \ 0--\qty{0.10}{s}: $\Phi$ rises from 0 to \qty{4.0}{mWb}; \
  \qtyrange{0.10}{0.25}{s}: steady; \ \qtyrange{0.25}{0.30}{s}: falls to 0}
\find{$\varepsilon$ in each part}
\formula{$\varepsilon = -N\dfrac{\Delta\Phi}{\Delta t}$, where
  $\dfrac{\Delta\Phi}{\Delta t}$ is the gradient}
\substitute{\eqlines{\text{0--0.10 s:}\ \frac{\Delta\Phi}{\Delta t}
    &= \frac{\qty{4.0e-3}{Wb} - 0}{\qty{0.10}{s}} = \qty{0.040}{Wb.s^{-1}}\\
  \varepsilon &= -50 \times \qty{0.040}{Wb.s^{-1}} = \qty{-2.0}{V}\\
  \text{0.10--0.25 s:}\ \frac{\Delta\Phi}{\Delta t} &= 0, \ \text{so}\ \varepsilon = 0\\
  \text{0.25--0.30 s:}\ \frac{\Delta\Phi}{\Delta t}
    &= \frac{0 - \qty{4.0e-3}{Wb}}{\qty{0.050}{s}} = \qty{-0.080}{Wb.s^{-1}}\\
  \varepsilon &= -50 \times (\qty{-0.080}{Wb.s^{-1}}) = +\qty{4.0}{V}}}
\answer{\final{$\qty{-2.0}{V}$, then 0, then $+\qty{4.0}{V}$}
  (\cref{ch3:fig:phigraph}(b))}
\checkstep{the last part is twice as steep, so its emf is twice as big and of the
  opposite sign \cmark. Each emf block's height $\times$ width is the same:
  $\qty{2.0}{V} \times \qty{0.10}{s} = \qty{4.0}{V} \times \qty{0.050}{s}
  = \qty{0.20}{V.s}$. This equals $N\lvert\Delta\Phi\rvert = 50 \times
  \qty{4.0e-3}{Wb} = \qty{0.20}{Wb}$, because $\qty{1}{Wb} = \qty{1}{V.s}$ \cmark.
  The signs only tell you that the emf in the last part is in the opposite
  direction to the first part. To find which way the current actually flows round
  the coil, use Lenz's law (\cref{ch3:sec:lenz}).}
\begin{figure}
  \centering
  \begin{tikzpicture}
    \begin{axis}[hsc, name=phi, width=0.47\linewidth, height=0.32\linewidth,
        xmin=0, xmax=0.35, ymin=0, ymax=5, xlabel={$t$ (\si{s})},
        ylabel={$\Phi$ (\si{mWb})}, xtick={0,0.1,0.2,0.3}, ytick={0,1,...,5},
        xticklabel style={/pgf/number format/.cd, fixed, precision=1},
        title={(a) flux}, title style={font=\small}]
      \addplot[hscseriesA, line width=1.3pt, mark=none]
        coordinates {(0,0) (0.10,4.0) (0.25,4.0) (0.30,0) (0.35,0)};
    \end{axis}
    \begin{axis}[hsc, at={(phi.east)}, anchor=west, xshift=12mm,
        width=0.47\linewidth, height=0.32\linewidth,
        xmin=0, xmax=0.35, ymin=-3, ymax=5, xlabel={$t$ (\si{s})},
        ylabel={$\varepsilon$ (\si{V})}, xtick={0,0.1,0.2,0.3}, ytick={-2,0,2,4},
        xticklabel style={/pgf/number format/.cd, fixed, precision=1},
        title={(b) induced emf}, title style={font=\small}]
      \draw[hscink!70, line width=0.6pt] (axis cs:0,0) -- (axis cs:0.35,0);
      \addplot[hscseriesB, solid, line width=1.3pt, mark=none, const plot]
        coordinates {(0,-2.0) (0.10,0) (0.25,4.0) (0.30,0) (0.35,0)};
    \end{axis}
  \end{tikzpicture}
  \caption{(a) Flux through each turn. (b) The emf is $-N \times$ the gradient
    of (a).}\label{ch3:fig:phigraph}
\end{figure}
\end{workedexample}

\begin{mistakes}
\begin{itemize}
  \item \xmark\ Using $\Phi$ instead of $\Delta\Phi$. \cmark\ Find
        $\Phi_\mathrm{initial}$ and $\Phi_\mathrm{final}$, then subtract.
  \item \xmark\ ``No emf, because the coil is not connected.'' \cmark\ The emf is
        still induced. Only a current needs a closed circuit.
  \item \xmark\ Using a clock reading as $\Delta t$. \cmark\ From
        $t = \qty{1.5}{s}$ to $t = \qty{4.0}{s}$, $\Delta t = \qty{2.5}{s}$.
\end{itemize}
\end{mistakes}

\begin{examtip}
Stating Faraday's law (1 mark): ``The induced emf is proportional to the rate of
change of magnetic flux through the circuit, $\varepsilon = -N\dfrac{\Delta\Phi}{\Delta t}$.''
In calculations (2--3 marks), show $\Phi_\mathrm{initial}$ and $\Phi_\mathrm{final}$,
substitute with units, and give the size with a unit. State the direction
separately, using Lenz's law.
\end{examtip}

\begin{quickcheck}
  \item The flux through each turn of a 120-turn coil falls from \qty{6.0}{mWb} to
        \qty{1.5}{mWb} in \qty{0.30}{s}. Find the size of the average emf.
  \item A 40-turn coil of area \qty{25}{cm^2} has its plane perpendicular to a
        \qty{0.15}{T} field. The field is switched off in \qty{0.050}{s}. Find the average emf,
        and the current if the total resistance is \qty{2.5}{\ohm}.
  \item A 500-turn coil must give an average emf of at least \qty{2.0}{V} when the
        flux through each turn changes by \qty{8.0e-4}{Wb}. Find the longest time
        the change can take.
  \item The ends of a coil are not joined. A magnet is pushed into it. Is there an
        induced emf? An induced current? Explain.
  \item The flux through a coil is very large but steady. Explain why the induced
        emf is zero.
\end{quickcheck}

%% ===================================================================
\section{Lenz's law: which way does the current flow?}\label{ch3:sec:lenz}
\covers{M6.IQ3.02}

\begin{bigidea}
An induced current always flows the way that fights the change that made it.
Its own magnetic field tries to keep the flux as it was. This is the minus sign
in Faraday's law, and it follows from conservation of energy.
\end{bigidea}

\begin{keywords}
  \kw{Lenz's law}{the induced current flows so that its magnetic field opposes
    the change in flux that produced it.}
  \kw{Induced magnetic field}{the magnetic field made by the induced current.}
  \kw{Viewed from \ldots}{the side you look from when you say ``clockwise''.}
\end{keywords}

Faraday's law gives the size of the emf. \term{Lenz's law} gives its direction.
Lenz's law has no formula of its own. It is the minus sign in Faraday's law
(\cref{ch3:sec:faraday}). Learn this sentence:
\begin{center}
  \begin{minipage}{0.9\linewidth}\itshape
  The induced current flows in the direction that makes a magnetic field
  opposing the change in flux that produced it.
  \end{minipage}
\end{center}
For an open circuit there is no current, so say: the induced emf acts in the
direction that would drive a current whose magnetic field opposes the change in
flux that produced it.

The induced field opposes the \emph{change}, not the field. If the flux is
increasing, the induced field points the opposite way to the original field. If
the flux is decreasing, it points the \emph{same} way, to ``top up'' the flux.

\begin{refresher}{The right-hand grip rule for a loop or coil}
\begin{sidebyside}[0.42]
  \centering
  \begin{tikzpicture}[x=1cm, y=1cm]
    \draw[line width=1.3pt, hscink, postaction={decorate}, decoration={markings,
      mark=at position 0.125 with {\arrow{Stealth[length=2.6mm, width=2mm]}},
      mark=at position 0.625 with {\arrow{Stealth[length=2.6mm, width=2mm]}}}]
      (0,0) circle (0.7);
    \foreach \p in {(0,0), (0.28,0.28), (-0.28,0.28), (0.28,-0.28), (-0.28,-0.28)}
      \fill[hscfield] \p circle (1.6pt);
    \node[font=\footnotesize, align=center] at (0,-1.15) {anticlockwise:\\ dots, N face};
    \begin{scope}[xshift=2.2cm]
      \draw[line width=1.3pt, hscink, postaction={decorate}, decoration={markings,
        mark=at position 0.125 with {\arrow{Stealth[length=2.6mm, width=2mm]}},
        mark=at position 0.625 with {\arrow{Stealth[length=2.6mm, width=2mm]}}}]
        (0.7,0) arc (0:-360:0.7);
      \foreach \p in {(0,0), (0.28,0.28), (-0.28,0.28), (0.28,-0.28), (-0.28,-0.28)}
        \node[text=hscfield, font=\small] at \p {$\times$};
      \node[font=\footnotesize, align=center] at (0,-1.15) {clockwise:\\ crosses, S face};
    \end{scope}
  \end{tikzpicture}
  \captionof{figure}{A loop seen from the front.}\label{ch3:fig:grip}
\nextside
  Curl the fingers of your right hand the way the current flows. Your thumb
  points along the field inside the loop.
  \begin{itemize}
    \item Anticlockwise (as you see it): field towards you. You face an N pole.
    \item Clockwise: field away from you. You face an S pole.
  \end{itemize}
\end{sidebyside}
\end{refresher}

\begin{sidebyside}[0.45]
  \centering
  \begin{tikzpicture}[x=0.8cm, y=0.8cm]
    \foreach \x in {-1.5,-0.9,-0.3,0.3,0.9,1.5} \foreach \y in {-1.5,-0.9,-0.3,0.3,0.9,1.5} {
      \pgfmathsetmacro\rr{abs(veclen(\x,\y)-1.1)}
      \ifdim\rr pt<0.2pt\else\fill[hscfield] (\x,\y) circle (1.6pt);\fi}
    \draw[line width=1.4pt, hscink, postaction={decorate}, decoration={markings,
      mark=at position 0.125 with {\arrow{Stealth[length=2.8mm, width=2.1mm]}},
      mark=at position 0.625 with {\arrow{Stealth[length=2.8mm, width=2.1mm]}}}]
      (0,0) circle (1.1);
    \node[font=\small, fill=white, inner sep=1pt] at (45:1.42) {$I$};
    \node[font=\small, fill=white, inner sep=1pt] at (225:1.42) {$I$};
    \draw[hscbrand, line width=0.9pt] (0,0) circle (5pt);
    \fill[hscbrand] (0,0) circle (1.8pt);
    \node[font=\small, align=center] at (0,1.95) {$B$ out of the page, decreasing};
  \end{tikzpicture}
  \captionof{figure}{Fewer dots, so the induced field (navy dot) points out of
    the page: anticlockwise current.}\label{ch3:fig:lenzloop}
\nextside
  \captionof{table}{Lenz's law: the induced field.}\label{ch3:tab:lenz}
  \centering
  \begin{tabular}{@{}L{26mm}L{40mm}@{}}
    \toprule
    \colhead{Flux change} & \colhead{Induced field inside the loop}\\
    \midrule
    more crosses & dots (out of the page)\\
    fewer crosses & crosses (into the page)\\
    more dots & crosses (into the page)\\
    fewer dots & dots (out of the page)\\
    \bottomrule
  \end{tabular}
\end{sidebyside}

Use three steps for every direction question. \Cref{ch3:fig:lenzloop} shows them for a loop.
\begin{enumerate}
  \item \textbf{Find the flux change}: which way does the field through the loop
        point, and is the flux increasing or decreasing?
  \item \textbf{Apply Lenz's law} to choose the induced field (\cref{ch3:tab:lenz}).
        Write the law down: it often earns a mark.
  \item \textbf{Use the grip rule} to find the current. Say where you view it from.
\end{enumerate}

For magnets, the \textbf{pole method} is quicker. The induced current turns the
coil into a weak magnet (\cref{ch3:fig:magnetcoil}). A magnet coming closer makes
the near face the \emph{same} pole (it repels). A magnet moving away makes the
near face the \emph{opposite} pole (it attracts). Either way the coil opposes the
motion. Then use \cref{ch3:fig:grip}: an N face has an anticlockwise current.

\begin{figure}
  \centering
  \begin{tikzpicture}[x=1cm, y=1cm]
    \foreach \xs in {0, 7.1} {
    \begin{scope}[xshift=\xs cm]
      \fill[hscvel!22] (0,0) rectangle (0.8,0.5);
      \fill[hscforce!22] (0.8,0) rectangle (1.6,0.5);
      \draw[line width=0.8pt, hscink] (0,0) rectangle (1.6,0.5);
      \node[font=\small\bfseries] at (0.4,0.25) {S};
      \node[font=\small\bfseries] at (1.2,0.25) {N};
      \foreach \x in {2.6, 2.85, 3.1, 3.35, 3.6}
        \draw[line width=1pt, hscink] (\x,0.25) ellipse (0.1 and 0.6);
    \end{scope}}
    % ----- (a) approaching
    \draw[velocity] (0.3,0.85) -- (1.3,0.85) node[right, text=hscvel] {$v$};
    \draw[->, hscfield, line width=1.2pt] (3.95,0.25) -- (2.05,0.25);
    \node[font=\footnotesize, text=hscfield, anchor=west] at (3.95,0.25) {induced $B$};
    \node[font=\footnotesize\bfseries, text=hscforce] at (3.1,1.15) {near face: N};
    \node[font=\footnotesize, align=center] at (2.2,-0.95)
      {(a) N pole approaching: the near\\ face becomes N and repels.};
    \chThreeACW{1.2}{-1.9}{0.36}
    \node[font=\footnotesize, align=left, anchor=west] at (1.75,-1.9)
      {seen from the magnet:\\ anticlockwise};
    % ----- (b) moving away
    \begin{scope}[xshift=7.1cm]
      \draw[velocity] (1.3,0.85) -- (0.3,0.85) node[left, text=hscvel] {$v$};
      \draw[->, hscfield, line width=1.2pt] (2.05,0.25) -- (3.95,0.25);
      \node[font=\footnotesize, text=hscfield, anchor=west] at (3.95,0.25) {induced $B$};
      \node[font=\footnotesize\bfseries, text=hscvel] at (3.1,1.15) {near face: S};
      \node[font=\footnotesize, align=center] at (2.2,-0.95)
        {(b) N pole moving away: the near\\ face becomes S and attracts.};
      \chThreeCW{1.2}{-1.9}{0.36}
      \node[font=\footnotesize, align=left, anchor=west] at (1.75,-1.9)
        {seen from the magnet:\\ clockwise};
    \end{scope}
  \end{tikzpicture}
  \caption{The pole method. The coil always opposes the motion of the
    magnet.}\label{ch3:fig:magnetcoil}
\end{figure}

\subsection{Lenz's law and conservation of energy}

Suppose Lenz's law worked the other way, so a coil \emph{attracted} a magnet
pushed towards it. The magnet would speed up, so the flux would change faster
and the current would grow. A bigger current pulls harder, and so on. The magnet
would gain kinetic energy and the coil would heat up, with no energy put in. That
breaks the \term{law of conservation of energy}: energy cannot be created or
destroyed, only changed in form.

So the coil must push back. You must do work to move the magnet, and that work
becomes electrical energy, then heat (\cref{ch3:fig:energy}). With an open
circuit there is no current, no push back, and no energy transfer.

\begin{figure}
  \centering
  \begin{tikzpicture}
    \node[concept=36mm] (a) {Work done by you\\ (against the opposing force)};
    \node[concept=34mm, right=9mm of a] (b) {Electrical energy\\ (induced current)};
    \node[concept=34mm, right=9mm of b] (c) {Heat in the wires\\ (or light in a lamp)};
    \draw[link] (a) -- (b);
    \draw[link] (b) -- (c);
  \end{tikzpicture}
  \caption{Energy transfers when a current is induced.}\label{ch3:fig:energy}
\end{figure}

\begin{workedexample}{Lifting a magnet away from a loop}\label{ch3:we:lift}
A wire loop lies flat on a table. A bar magnet is held above it, S pole down,
and is lifted straight up. Find the direction of the induced current, viewed
from above.
\given{field lines enter an S pole, so under the magnet the field points
  \emph{up} through the loop. Lifting the magnet weakens it.}
\find{the current direction, viewed from above}
\formula{Lenz's law, then the right-hand grip rule}
\substitute{Step 1: the upward flux through the loop is decreasing.\\
  Step 2: by Lenz's law the induced current opposes this change. The flux is
  decreasing, so the induced field points the same way: up.\\
  Step 3: an upward field comes towards a viewer above, so the grip rule gives
  an anticlockwise current.}
\answer{\final{anticlockwise, viewed from above}}
\checkstep{pole method: anticlockwise means the top face is N. It attracts the S
  pole and pulls back on the rising magnet, opposing the motion \cmark.}
\end{workedexample}

\begin{workedexample}{Pulling a loop out of a field}\label{ch3:we:pullout}
A rectangular loop is pulled to the right at a steady speed. Its left half is
still in a field that points into the page (\cref{ch3:fig:pullout}). Find (a) the
direction of the induced current and (b) the direction of the magnetic force on
the loop. (c) Explain the energy transfers.
\begin{figure}
  \centering
  \begin{tikzpicture}[x=1cm, y=1cm]
    \draw[hscfield, dashed, line width=0.6pt] (0,0) rectangle (2.6,2.0);
    \foreach \x in {0.3, 0.9, 1.5, 2.2} \foreach \y in {0.3, 0.75, 1.25, 1.7}
      \node[text=hscfield, font=\small] at (\x,\y) {$\times$};
    \draw[line width=1.6pt, hscbrand] (1.7,0.55) rectangle (3.5,1.45);
    \draw[velocity] (3.5,1.0) -- (4.5,1.0) node[right, text=hscvel] {$v$ (pulled)};
    \node[font=\small, anchor=west] at (-3.4,1.0) {$B$ into the page};
  \end{tikzpicture}
  \caption{The loop is leaving the field region (dashed).}\label{ch3:fig:pullout}
\end{figure}
\given{$B$ into the page; the area of the loop inside the field is shrinking.}
\find{(a) current direction; (b) force direction; (c) energy changes}
\formula{$\Phi = BA\cos\theta$; Lenz's law; grip rule; $F = lIB\sin\theta$ with
  the right-hand palm rule}
\substitute{(a) Less of the loop is in the field, so the into-page flux
  decreases. By Lenz's law the induced current opposes this decrease, so its
  field points into the page inside the loop. The grip rule then gives a
  clockwise current.\\
  (b) Only the left side is fully in the field. Clockwise means its current goes
  up. Palm rule: fingers into the page, thumb up, palm pushes left. The forces on
  the top and bottom sides are equal and opposite, so they cancel.\\
  (c) The pull must do work against this force. The work becomes electrical
  energy, then heat in the wire.}
\answer{(a) \final{clockwise} \quad (b) \final{to the left} \quad (c) work
  $\to$ electrical energy $\to$ heat}
\checkstep{the cause is the pull to the right, and the force opposes it \cmark.}
\end{workedexample}

\begin{mistakes}
\begin{itemize}
  \item \xmark\ ``The induced field always opposes the magnet's field.''
        \cmark\ It opposes the \emph{change}; for a decreasing flux it points the
        same way.
  \item \xmark\ ``Clockwise'' with no viewpoint. \cmark\ ``Clockwise, viewed from
        above.''
  \item \xmark\ Treating an S pole like an N pole. \cmark\ Field lines point
        \emph{into} an S pole.
  \item \xmark\ Using your left hand. \cmark\ The grip and palm rules use the
        right hand.
\end{itemize}
\end{mistakes}

\begin{examtip}
``Account for Lenz's law using conservation of energy'' (3 marks) needs three
linked points: (1) state Lenz's law; (2) because the current opposes the change,
work must be done to cause it, and this work becomes electrical energy and heat;
(3) if the current helped the change, the motion would speed itself up and make
energy from nothing, which breaks conservation of energy.
In multiple choice, ``opposes the relative motion'' is true for moving magnets.
The full statement is ``opposes the \emph{change in flux}''. It also covers two
coils where nothing moves.
\end{examtip}

\begin{quickcheck}[breakable]
  \item State Lenz's law.
  \item A long straight wire lies in the page and carries a current to the right.
        A square loop lies in the page just above the wire. The loop moves up the
        page, away from the wire. Find the direction of the induced current.
        (Hint: use the right-hand grip rule for a straight wire
        (\cref{ch2:sec:fields}): thumb along the current, fingers curl the way the
        field goes. Also, $B = \dfrac{\mu_0 I}{2\pi r}$ gets smaller as $r$
        grows.)\label{ch3:qc:wire}
  \item The S pole of a magnet is pulled away from a coil. Which pole does the near
        face of the coil become? Which way does the current flow, viewed from the
        magnet?
  \item A loop lies in the page, in a field into the page. The loop is squashed,
        so its area gets smaller. Find the direction of the induced current.
  \item A student claims she can design a coil whose induced current
        \emph{attracts} a magnet that is pushed towards it. Use the law of
        conservation of energy to explain why this is impossible. (3 marks)
\end{quickcheck}

%% ===================================================================
\section{Motional emf: a straight wire moving in a field}\label{ch3:sec:motional}
\covers{M6.IQ3.02}

\begin{bigidea}
A metal rod moving across a magnetic field has an emf between its ends. The
field pushes the free charges (charges that can move through the metal, mainly
electrons) along the rod. In a closed circuit this drives a
current, and the field then pushes back on the rod: Lenz's law again.
\end{bigidea}

\begin{keywords}
  \kw{Motional emf}{the emf in a conductor because it moves across a magnetic
    field.}
  \kw{Conventional current}{the direction positive charge would flow. Electrons
    flow the other way.}
\end{keywords}

\begin{refresher}{Force on a moving charge (\cref{ch1:chap,ch2:chap})}
A charge moving across a field feels $F = qv_\perp B = qvB\sin\theta$.
\term{Right-hand palm rule}: point your fingers along $B$ and your thumb along
the velocity of a \emph{positive} charge (or the current). Your palm then pushes
in the direction of the force. The force on an electron is the opposite way.
\end{refresher}

A metal rod moves right through a field into the page (\cref{ch3:fig:rod}(a)).
\begin{itemize}
  \item The charges move right with the rod. Palm rule for a positive charge:
        fingers into the page, thumb right, palm up. Positive charge is pushed
        \emph{up}; the free electrons (the electrons that can move through the
        metal) are pushed down.
  \item The bottom end becomes $-$ and the top end $+$. The rod is now like a
        battery: this is \term{motional emf}.
  \item The build-up stops when the electric field of the separated charges
        balances the magnetic force.
\end{itemize}

\begin{figure}
  \centering
  \begin{tikzpicture}[x=1cm, y=1cm]
    % ---------- (a) rod on its own
    \begin{scope}
      \foreach \x in {0.3, 0.9, 2.1, 2.7, 3.3} \foreach \y in {0.3, 0.9, 1.5, 2.1, 2.7}
        \node[text=hscfield, font=\small] at (\x,\y) {$\times$};
      \draw[line width=3.2pt, hscmuted] (1.5,0.15) -- (1.5,2.85);
      \node[draw, circle, inner sep=0.6pt, font=\footnotesize\bfseries, fill=white]
        at (1.5,3.15) {$+$};
      \node[draw, circle, inner sep=0.6pt, font=\footnotesize\bfseries, fill=white]
        at (1.5,-0.15) {$-$};
      \draw[velocity] (1.62,1.2) -- (2.5,1.2);
      \node[text=hscvel, font=\small, fill=white, inner sep=1pt, anchor=west]
        at (2.52,1.2) {$v$};
      \draw[force] (1.2,1.25) -- (1.2,2.35);
      \node[font=\footnotesize, text=hscforce, fill=white, inner sep=1pt, anchor=east]
        at (1.12,1.8) {$F$ on $+q$};
      \node[align=center, text width=4.0cm, font=\small, anchor=north] at (1.8,-1.1)
        {(a) Rod alone: the top end becomes $+$.};
    \end{scope}
    % ---------- (b) rod on rails
    \begin{scope}[xshift=5.6cm]
      \fill[hscvel!12] (3.0,0) rectangle (3.8,3.0);
      \foreach \x in {0.9, 1.65, 2.4, 4.6, 5.4} \foreach \y in {0.5, 1.5, 2.5}
        \node[text=hscfield, font=\small] at (\x,\y) {$\times$};
      \foreach \y in {0.5, 2.5} \node[text=hscfield, font=\small] at (3.4,\y) {$\times$};
      \draw[line width=1.5pt, hscink] (0,0) -- (6.0,0);
      \draw[line width=1.5pt, hscink] (0,3.0) -- (6.0,3.0);
      \draw[line width=1.2pt, hscink] (0,0) -- (0,1.0);
      \draw[line width=1.2pt, hscink] (0,2.0) -- (0,3.0);
      \draw[line width=1.2pt, hscink, fill=white] (-0.16,1.0) rectangle (0.16,2.0);
      \node[font=\small, anchor=west] at (0.2,1.5) {$R$};
      \draw[line width=1pt, hscmuted, dashed] (3.0,-0.2) -- (3.0,3.2);
      \draw[line width=3.2pt, hscmuted] (3.8,-0.2) -- (3.8,3.2);
      \draw[velocity] (3.92,2.0) -- (4.9,2.0);
      \node[text=hscvel, font=\small, fill=white, inner sep=1pt, anchor=west]
        at (4.92,2.0) {$v$};
      \draw[force] (3.68,1.0) -- (2.7,1.0);
      \node[font=\footnotesize, text=hscforce, fill=white, inner sep=1pt, anchor=east]
        at (2.65,1.0) {$F = lIB$};
      \draw[->, line width=1pt, hscink] (4.05,0.2) -- (4.05,0.95)
        node[right, font=\small] {$I$};
      \draw[->, line width=1pt, hscink] (2.2,3.28) -- (1.2,3.28) node[left, font=\small] {$I$};
      \draw[->, line width=1pt, hscink] (1.2,-0.28) -- (2.2,-0.28) node[right, font=\small] {$I$};
      \draw[->, line width=1pt, hscink] (-0.45,2.0) -- (-0.45,1.0)
        node[midway, left, font=\small] {$I$};
      \draw[<->, line width=0.6pt, hscmuted] (6.35,0) -- (6.35,3.0)
        node[midway, right, text=hscink, font=\small] {$l$};
      \draw[<->, line width=0.6pt, hscmuted] (3.0,-0.62) -- (3.8,-0.62);
      \node[font=\small, below] at (3.4,-0.64) {$v\Delta t$};
      \node[align=center, text width=6.6cm, font=\small, anchor=north] at (3.0,-1.1)
        {(b) On rails: the current is anticlockwise. The rod sweeps the shaded
        area in time $\Delta t$.};
    \end{scope}
  \end{tikzpicture}
  \caption{Motional emf. The field points into the page.}\label{ch3:fig:rod}
\end{figure}

Now put the rod on metal rails joined by a resistor (\cref{ch3:fig:rod}(b)). A
current flows up the rod and down through the resistor.
\begin{itemize}
  \item The current in the rod feels a motor-effect force, $F = lIB$. Palm rule:
        fingers into the page, thumb up, palm left. The force \emph{opposes} the
        motion (Lenz). Something must keep pushing the rod, and that work becomes
        heat in the resistor.
  \item Faraday's view: the rod sweeps out area, so the flux through the circuit
        grows. More crosses means an anticlockwise current. Both views agree.
\end{itemize}

\begin{formulacard}[notonsheet]{Motional emf}
  \cardformula{\varepsilon = lvB}
  \begin{symbols}
    \sym{$\varepsilon$}{emf between the ends of the rod}{\si{V}}
    \sym{$l$}{length of the rod inside the field}{\si{m}}
    \sym{$v$}{speed of the rod across the field}{\si{m.s^{-1}}}
    \sym{$B$}{magnetic field strength}{\si{T}}
  \end{symbols}
  \inwords{a longer rod, moving faster through a stronger field, gives a bigger
    emf.}
  \whentouse{a straight rod in a uniform field, with the rod, its velocity and the
    field at \ang{90} to each other.}
  \whennottouse{the rod moves along the field lines or along its own length: then
    $\varepsilon = 0$. At an angle, use only the part of $v$ that is perpendicular
    to both $B$ and the rod.}
  \rearranged{$v = \dfrac{\varepsilon}{lB}$ \qquad $B = \dfrac{\varepsilon}{lv}$
    \qquad $l = \dfrac{\varepsilon}{vB}$}
  \tensecond{$l = \qty{0.50}{m}$, $v = \qty{4.0}{m.s^{-1}}$, $B = \qty{0.20}{T}$: \
    $\varepsilon = 0.50 \times 4.0 \times 0.20 = \qty{0.40}{V}$.}
  \derive{Faraday's law for the one-turn circuit in \cref{ch3:fig:rod}(b). In time
    $\Delta t$ the rod moves $v\Delta t$ and sweeps an area $\Delta A = l\,v\Delta t$.
    \begin{align*}
      \Delta\Phi &= B\,\Delta A = B\,l\,v\Delta t && \text{only $A$ changes; $\theta = \ang{0}$}\\
      \lvert\varepsilon\rvert &= N\frac{\Delta\Phi}{\Delta t} = \frac{B\,l\,v\Delta t}{\Delta t}
        && \text{Faraday's law, $N = 1$}\\
      \varepsilon &= lvB && \text{cancel $\Delta t$}
    \end{align*}}
  \trap{$l$ is the length \emph{inside the field}, not the whole wire. Use the palm
    rule for positive charges.}
\end{formulacard}

\begin{workedexample}{A rod moving through a field}\label{ch3:we:rod}
A \qty{0.40}{m} metal rod moves right at \qty{3.0}{m.s^{-1}} through a
\qty{0.25}{T} field into the page, as in \cref{ch3:fig:rod}(a). Find the emf
between its ends, and which end is positive.
\given{$l = \qty{0.40}{m}$, \ $v = \qty{3.0}{m.s^{-1}}$, \ $B = \qty{0.25}{T}$, all at
  \ang{90} to each other}
\find{$\varepsilon$ and the positive end}
\formula{$\varepsilon = lvB$; right-hand palm rule}
\substitute{$\varepsilon = \qty{0.40}{m} \times \qty{3.0}{m.s^{-1}} \times \qty{0.25}{T}
  = \qty{0.30}{V}$.\\
  Palm rule on a positive charge: fingers into the page, thumb right, palm up.}
\answer{\final{$\varepsilon = \qty{0.30}{V}$}, top end positive}
\checkstep{$\unit{m} \times \unit{m.s^{-1}} \times \unit{T} = \unit{Wb.s^{-1}} =
  \unit{V}$ \cmark. With no circuit there is no current and no opposing force.}
\end{workedexample}

\begin{workedexample}{A rod on rails: forces and energy}\label{ch3:we:rails}
A \qty{25}{cm} rod slides right at a steady \qty{2.0}{m.s^{-1}} on rails joined by
a \qty{0.50}{\ohm} resistor (\cref{ch3:fig:rod}(b)). $B = \qty{0.80}{T}$ into the
page; ignore friction and all other resistance. Find (a) the emf, (b) the current,
(c) the force needed to keep the rod moving, and (d) show that the power of this
force equals the electrical power.
\given{$l = \qty{0.25}{m}$, \ $v = \qty{2.0}{m.s^{-1}}$, \ $B = \qty{0.80}{T}$, \
  $R = \qty{0.50}{\ohm}$}
\find{$\varepsilon$, $I$, the push $F$, and the two powers}
\formula{$\varepsilon = lvB$, \ $I = \dfrac{\varepsilon}{R}$, \
  $F = lIB\sin\theta$ ($\theta = \ang{90}$),\\
  $P = F_\parallel v = Fv\cos\theta$ ($\theta = \ang{0}$: push along $v$), \
  $P = VI$}
\substitute{\eqlines{\varepsilon &= \qty{0.25}{m} \times \qty{2.0}{m.s^{-1}} \times
    \qty{0.80}{T} = \qty{0.40}{V}\\
  I &= \frac{\qty{0.40}{V}}{\qty{0.50}{\ohm}} = \qty{0.80}{A}\\
  F &= \qty{0.25}{m} \times \qty{0.80}{A} \times \qty{0.80}{T} = \qty{0.16}{N}
    \ \text{(magnetic force, to the left)}\\
  F_\mathrm{push} &= \qty{0.16}{N}\ \text{to the right (steady speed, so the forces
    balance)}\\
  P_\mathrm{push} &= F_\mathrm{push}v = \qty{0.16}{N} \times \qty{2.0}{m.s^{-1}}
    = \qty{0.32}{W}\\
  P_\mathrm{elec} &= VI = \qty{0.40}{V} \times \qty{0.80}{A} = \qty{0.32}{W}}}
\answer{(a) \final{\qty{0.40}{V}} (b) \final{\qty{0.80}{A}}, anticlockwise
  (c) a push of \final{\qty{0.16}{N} to the right} (d) both \qty{0.32}{W}}
\checkstep{at a steady speed the push (to the right) balances the magnetic force
  on the rod (to the left). Every joule of
  work done by the push becomes electrical energy, then heat: energy is
  conserved \cmark.}
\end{workedexample}

\begin{mistakes}
\begin{itemize}
  \item \xmark\ Using the palm rule with the electrons' velocity. \cmark\ Use
        positive charges, or reverse the answer.
  \item \xmark\ Using the whole length of the wire. \cmark\ Only the length in the
        field counts.
  \item \xmark\ ``A straight wire has no area, so this is not Faraday's law.''
        \cmark\ The rod is part of a circuit whose area changes; $\varepsilon = lvB$
        comes from Faraday's law.
\end{itemize}
\end{mistakes}

\begin{examtip}
Direction questions (1 mark) want the palm rule on positive charges and a clear
answer (``current flows up the rod''). A common 4-mark question compares two
rods. Rod 1 is pushed through a field. Rod 2 carries a current from a battery and
starts to move. Full marks need four points:
\begin{enumerate}
  \item rod 1: the field pushes the moving charges along the rod, so an emf is
        induced. The work done by the push (or the rod's kinetic energy, if it
        slows down) $\to$ electrical energy;
  \item rod 2: the motor effect, $F = lIB$, pushes the rod (electrical
        $\to$ kinetic energy);
  \item one similarity: both are forces on moving charges in a magnetic field;
  \item one difference: in rod 1 motion causes a current; in rod 2 a current
        causes motion.
\end{enumerate}
$\varepsilon = lvB$ is \emph{not} on the sheet, so show where it comes from.
\end{examtip}

\begin{quickcheck}
  \item An aircraft with a \qty{32}{m} wingspan flies horizontally at
        \qty{240}{m.s^{-1}}. The vertical part of Earth's field is \qty{5.0e-5}{T}.
        Find the emf between the wing tips.
  \item A rod moves parallel to the field lines. What is the emf? Explain.
  \item A horizontal rod moves up the page through a field that points out of the
        page. Which end becomes positive?
  \item A \qty{0.50}{m} rod slides along frictionless rails. A \qty{0.40}{T}
        field is at \ang{90} to the rails and the rod. The rails are joined by a \qty{2.0}{\ohm} resistor; ignore all other
        resistance. Find the speed that gives a current of \qty{0.15}{A}. Then
        find the push needed to keep the rod moving at this steady speed.
\end{quickcheck}

%% ===================================================================
\section{Solenoids, metal plates and two coils}\label{ch3:sec:solenoids}
\covers{M6.IQ3.02}

\begin{bigidea}
Any change in flux induces an emf, whatever causes it. A magnet can pass through
a solenoid. A metal plate can move through a field. The current in a nearby coil
can change. All three follow Faraday's and Lenz's laws.
\end{bigidea}

\begin{keywords}
  \kw{Solenoid}{a long coil of wire.}
  \kw{Eddy currents}{loops of current induced inside a solid piece of metal when
    the flux through it changes.}
  \kw{Mutual induction}{an emf induced in one coil (the secondary) by a changing
    current in a nearby coil (the primary).}
\end{keywords}

\subsection{A magnet moving through a solenoid}

As a magnet comes towards a coil, the flux through the coil increases. This
induces an emf in one direction. When the magnet is at the middle, the flux is at
its largest. For an instant it is not changing, so the emf is zero. As the magnet
leaves, the flux decreases. This induces an emf in the other direction. So a
magnet passing through a coil gives two pulses of opposite sign
(\cref{ch3:fig:falling}, \cref{ch3:we:drop}). A faster magnet gives taller,
narrower pulses.

Exam questions often draw the coil from the side. They ask which way the current
flows through a meter. Work in two stages. First, use Lenz's law to find the
current as seen from above. Second, change this into a direction along the wires
you can see. The \emph{front} of each turn is the half nearest you. Seen from
above, a clockwise current moves from right to left along the front
(\cref{ch3:fig:front}). Follow it along the front wires to the leads
(\cref{ch3:we:winding}).

\begin{figure}
  \centering
  \begin{tikzpicture}[x=1cm, y=1cm]
    % ---------- (a) top view of one turn, clockwise
    \begin{scope}
      \draw[line width=0.8pt, hscmuted!70] (0.9,0) arc (0:180:0.9);
      \draw[line width=1.6pt, hscink] (-0.9,0) arc (180:360:0.9);
      \draw[line width=0.8pt, hscmuted!70, -{Stealth[length=2.4mm, width=2mm]}]
        (-0.1,0.895) -- (0.05,0.9);
      \draw[line width=1.6pt, hscink, -{Stealth[length=2.8mm, width=2.2mm]}]
        (0.1,-0.895) -- (-0.05,-0.9);
      \node[font=\footnotesize, text=hscmuted, anchor=south] at (0,0.95) {back};
      \node[font=\footnotesize, anchor=west] at (0.8,-0.65) {front};
      \draw[->, line width=0.9pt, hscink] (0,-1.85) -- (0,-1.2);
      \node[font=\footnotesize, anchor=north] at (0,-1.85) {you look this way};
      \node[font=\small, align=center, anchor=north] at (0,-2.3)
        {(a) Top view: clockwise};
    \end{scope}
    % ---------- (b) the same turn seen from the side (slightly above)
    \begin{scope}[xshift=5.6cm, yshift=-0.3cm]
      \draw[line width=0.8pt, hscmuted!70, dashed] (-1.4,0) arc (180:0:1.4 and 0.42);
      \draw[line width=1.6pt, hscink] (1.4,0) arc (0:-180:1.4 and 0.42);
      \draw[line width=1.6pt, hscink, -{Stealth[length=2.8mm, width=2.2mm]}]
        (0.1,-0.419) -- (-0.05,-0.42);
      \draw[line width=0.8pt, hscmuted!70, -{Stealth[length=2.4mm, width=2mm]}]
        (-0.1,0.419) -- (0.05,0.42);
      \node[font=\footnotesize, text=hscmuted, anchor=south] at (0,0.47) {back};
      \node[font=\footnotesize, anchor=north, align=center] at (0,-0.52)
        {front: right to left};
      \node[font=\small, align=center, anchor=north] at (0,-2.0)
        {(b) Side view: what you see};
    \end{scope}
  \end{tikzpicture}
  \caption{One turn of a coil carrying a clockwise current (seen from above).
    Along the front, the half nearest you, the current flows from right to
    left.}\label{ch3:fig:front}
\end{figure}

\subsection{Metal plates: eddy currents}

A metal plate is full of possible current loops. When the flux through part of
the plate changes, an emf is induced there. Currents then swirl round inside the
metal. These are \term{eddy currents}: loops of current inside solid metal. They
are named after eddies, the small whirlpools in a river (\cref{ch3:fig:plate}).

Eddy currents form only where the flux through the metal is changing. Use Lenz's
law to find their direction. In \cref{ch3:fig:plate}(a) the plate is entering
the field. The into-page flux through its right-hand part increases. The induced
field points out of the page, so the eddy current is anticlockwise. In (b) the
plate is fully inside and the flux does not change, so there are no eddy
currents. In (c) the plate is leaving the field. The into-page flux through its
left-hand part decreases, so the induced field points into the page and the eddy
current is clockwise.

The part of each eddy loop that is inside the field feels a motor-effect force.
This force points backwards, against the motion. So the plate slows down, and its
kinetic energy becomes heat in the metal. This is \textbf{magnetic braking}
(\cref{ch5:chap}).

\begin{figure}
  \centering
  \begin{tikzpicture}[x=1cm, y=1cm]
    % ---------- (a) entering: field region on the right of x = 1.6
    \begin{scope}
      \draw[hscfield, dashed, line width=0.6pt] (1.6,0) -- (1.6,2.3) -- (4.2,2.3);
      \draw[hscfield, dashed, line width=0.6pt] (1.6,0) -- (4.2,0);
      \fill[hscmuted!22] (0.7,0.4) rectangle (2.5,1.9);
      \draw[line width=0.9pt, hscink] (0.7,0.4) rectangle (2.5,1.9);
      \foreach \x in {1.85, 2.3, ..., 4.0} \foreach \y in {0.25, 0.7, 1.15, 1.6, 2.05}
        \node[text=hscfield, font=\small] at (\x,\y) {$\times$};
      \chThreeACW{1.6}{1.15}{0.45}
      \draw[velocity] (1.15,2.6) -- (2.05,2.6) node[right, text=hscvel] {$v$};
      \draw[force] (2.1,-0.22) -- (1.1,-0.22) node[left, text=hscforce] {$F$};
      \node[align=center, text width=4.1cm, font=\footnotesize] at (2.1,-0.85)
        {(a) Entering: flux increasing.\\ Anticlockwise eddy current.};
    \end{scope}
    % ---------- (b) fully inside
    \begin{scope}[xshift=4.85cm]
      \draw[hscfield, dashed, line width=0.6pt] (0,0) rectangle (4.2,2.3);
      \fill[hscmuted!22] (1.2,0.4) rectangle (3.0,1.9);
      \draw[line width=0.9pt, hscink] (1.2,0.4) rectangle (3.0,1.9);
      \foreach \x in {0.25, 0.7, ..., 4.0} \foreach \y in {0.25, 0.7, 1.15, 1.6, 2.05}
        \node[text=hscfield, font=\small] at (\x,\y) {$\times$};
      \draw[velocity] (1.65,2.6) -- (2.55,2.6) node[right, text=hscvel] {$v$};
      \node[align=center, text width=4.1cm, font=\footnotesize] at (2.1,-0.85)
        {(b) Fully inside: flux not changing.\\ No eddy current, no force.};
    \end{scope}
    % ---------- (c) leaving: field region on the left of x = 2.6
    \begin{scope}[xshift=9.7cm]
      \draw[hscfield, dashed, line width=0.6pt] (0,0) -- (2.6,0) -- (2.6,2.3) -- (0,2.3);
      \fill[hscmuted!22] (1.7,0.4) rectangle (3.5,1.9);
      \draw[line width=0.9pt, hscink] (1.7,0.4) rectangle (3.5,1.9);
      \foreach \x in {0.25, 0.7, ..., 2.1} \foreach \y in {0.25, 0.7, 1.15, 1.6, 2.05}
        \node[text=hscfield, font=\small] at (\x,\y) {$\times$};
      \chThreeCW{2.6}{1.15}{0.45}
      \draw[velocity] (2.15,2.6) -- (3.05,2.6) node[right, text=hscvel] {$v$};
      \draw[force] (3.1,-0.22) -- (2.1,-0.22) node[left, text=hscforce] {$F$};
      \node[align=center, text width=4.1cm, font=\footnotesize] at (2.1,-0.85)
        {(c) Leaving: flux decreasing.\\ Clockwise eddy current.};
    \end{scope}
  \end{tikzpicture}
  \caption{A metal plate moving right past the edges of a field region (dashed,
    $B$ into the page). Eddy currents form only where the flux through the metal
    changes. The force $F$ on them opposes the motion.}\label{ch3:fig:plate}
\end{figure}

Evidence for Lenz's law: a solid copper plate swinging between magnet poles stops
in a few swings. A plate with slots cut in it swings much longer, because the
slots break the current loops. A strong magnet falls slowly through a copper pipe
(\cref{ch3:sec:experiments}). Eddy currents are useful in brakes, induction
cooktops and metal detectors, but waste energy in transformer cores (\cref{ch4:chap}).

\subsection{One solenoid near another}

Put two coils on one iron core (\cref{ch3:fig:twocoils}). The \term{primary coil}
is the coil joined to the supply; here it has a battery and a switch. The
\term{secondary coil} is the coil where the emf is induced; here it has a
galvanometer. Inducing an emf in one coil with a changing current in another is
called \term{mutual induction}.
\begin{itemize}
  \item \textbf{Switch closed:} the primary current and its field grow, so the flux
        through the secondary increases. A brief emf is induced.
  \item \textbf{Kept closed:} the current is steady, so the flux is steady. There
        is no emf.
  \item \textbf{Switch opened:} the flux decreases. A brief emf is induced the
        other way.
\end{itemize}

By Lenz's law, while the primary's field grows, the secondary's induced field
points the opposite way. While the primary's field falls, the induced field
points the same way. The core is made of \term{soft iron}: iron that is easy to
magnetise and loses its magnetism when the current stops. It guides the field
lines through the secondary. So more of the primary's flux passes through the
secondary, and the emf is bigger.

Two more ways to change the flux through the secondary:
\begin{itemize}
  \item \textbf{Moving one coil:} use two separate coils. Keep a steady current in
        a small primary coil and push it into a larger secondary coil. The flux
        through the secondary increases, so a brief emf is induced. Hold it
        still: no emf. Pull it out: a brief emf the other way. Faster motion
        gives a bigger emf. The primary acts just like a bar magnet, so the pole
        method (\cref{ch3:fig:magnetcoil}) gives the direction.
  \item \textbf{AC supply:} an \term{AC} (alternating current) supply makes a
        current that keeps changing size and reversing direction. A \term{DC}
        (direct current) supply, like a battery, makes a steady current in one
        direction. An AC primary current gives a changing flux all the time, so
        the secondary has an emf all the time (an AC emf). This is how a
        transformer works (\cref{ch4:chap}).
\end{itemize}

\textbf{Energy:} the electrical energy in the secondary circuit comes from the
primary's power supply. (If you move one coil, it comes from the work you do.) While a
current flows in the secondary, its field opposes the change in flux. So the
primary supply must do extra work to keep its current changing. Energy is
transferred from the primary circuit to the secondary circuit, then becomes heat
in the secondary. No energy is created.

\begin{figure}
  \centering
  \begin{tikzpicture}[x=1cm, y=1cm]
    \fill[hscmuted!30] (0,0) rectangle (6.4,0.6);
    \draw[line width=0.7pt, hscink] (0,0) rectangle (6.4,0.6);
    \node[font=\footnotesize, anchor=south] at (3.2,0.62) {soft iron core};
    \foreach \x in {0.5, 0.75, ..., 2.0}
      \draw[line width=1.1pt, hscink] (\x,-0.12) -- (\x+0.15,0.72);
    \foreach \x in {4.25, 4.5, ..., 5.75}
      \draw[line width=1.1pt, hscink] (\x,-0.12) -- (\x+0.15,0.72);
    \node[font=\footnotesize, anchor=south] at (1.3,0.95) {primary};
    \node[font=\footnotesize, anchor=south] at (5.05,0.95) {secondary};
    % primary circuit: battery and switch
    \draw[line width=0.9pt, hscink] (0.5,-0.12) -- (0.5,-1.1) -- (0.9,-1.1);
    \draw[line width=1.0pt, hscink] (0.9,-0.8) -- (0.9,-1.4);
    \draw[line width=2.4pt, hscink] (1.1,-0.95) -- (1.1,-1.25);
    \draw[line width=0.9pt, hscink] (1.1,-1.1) -- (1.4,-1.1);
    \fill[hscink] (1.4,-1.1) circle (1.4pt);
    \draw[line width=0.9pt, hscink] (1.4,-1.1) -- ++(25:0.56);
    \fill[hscink] (2.0,-1.1) circle (1.4pt);
    \draw[line width=0.9pt, hscink] (2.0,-1.1) -- (2.0,-0.12);
    \node[font=\footnotesize] at (1.0,-1.7) {battery};
    \node[font=\footnotesize] at (1.45,-0.55) {switch};
    % secondary circuit: galvanometer
    \draw[line width=0.9pt, hscink] (4.25,-0.12) -- (4.25,-1.1) -- (4.72,-1.1);
    \draw[line width=0.9pt, hscink] (5.38,-1.1) -- (5.9,-1.1) -- (5.9,-0.12);
    \draw[line width=0.9pt, hscink, fill=white] (5.05,-1.1) circle (0.33);
    \node[font=\small] at (5.05,-1.1) {G};
    % graphs
    \begin{scope}[xshift=7.5cm, yshift=-1.6cm]
    \begin{axis}[hsc, name=ip, width=6.2cm, height=2.6cm, xmin=0, xmax=6, ymin=-0.2, ymax=1.3,
        xtick=\empty, ytick=\empty, ylabel={$I_\mathrm{p}$},
        ylabel style={rotate=-90, font=\small}]
      \addplot[hscseriesA, line width=1.2pt, mark=none] coordinates {(0,0) (1,0)};
      \addplot[hscseriesA, line width=1.2pt, mark=none, domain=1:4, samples=60] {1-exp(-(x-1)/0.12)};
      \addplot[hscseriesA, line width=1.2pt, mark=none, domain=4:6, samples=60] {exp(-(x-4)/0.06)};
      \draw[guide] (axis cs:1,-0.2) -- (axis cs:1,1.3);
      \draw[guide] (axis cs:4,-0.2) -- (axis cs:4,1.3);
      \node[font=\footnotesize, anchor=south] at (axis cs:1,1.3) {closed};
      \node[font=\footnotesize, anchor=south] at (axis cs:4,1.3) {opened};
    \end{axis}
    \begin{axis}[hsc, at={(ip.south)}, anchor=north, yshift=-3mm, width=6.2cm, height=3.0cm,
        xmin=0, xmax=6, ymin=-1.3, ymax=2.3, xtick=\empty, ytick=\empty,
        xlabel={time $t$}, ylabel={$\varepsilon_\mathrm{s}$},
        ylabel style={rotate=-90, font=\small}]
      \draw[hscink!70, line width=0.6pt] (axis cs:0,0) -- (axis cs:6,0);
      \addplot[hscseriesB, solid, line width=1.2pt, mark=none] coordinates {(0,0) (1,0) (1,-1)};
      \addplot[hscseriesB, solid, line width=1.2pt, mark=none, domain=1:4, samples=120] {-exp(-(x-1)/0.12)};
      \addplot[hscseriesB, solid, line width=1.2pt, mark=none] coordinates {(4,0) (4,2)};
      \addplot[hscseriesB, solid, line width=1.2pt, mark=none, domain=4:6, samples=120] {2*exp(-(x-4)/0.06)};
      \draw[guide] (axis cs:1,-1.3) -- (axis cs:1,2.3);
      \draw[guide] (axis cs:4,-1.3) -- (axis cs:4,2.3);
      \node[font=\footnotesize, anchor=south] at (axis cs:2.5,0.1) {steady: zero};
    \end{axis}
    \end{scope}
  \end{tikzpicture}
  \caption{Two coils on one iron core. Right: the primary current $I_\mathrm{p}$
    and the secondary emf $\varepsilon_\mathrm{s}$ against time. The secondary has
    an emf only while the primary current changes. Opening the switch usually stops
    the current faster, so that pulse is taller and narrower.}\label{ch3:fig:twocoils}
\end{figure}

\begin{refresher}{The field inside a solenoid (on the formula sheet)}
Inside a long solenoid with an air core, $B = \dfrac{\mu_0 NI}{L}$. Here $N$ is the
number of turns of \emph{that} solenoid, $I$ its current, $L$ its length (\si{m}),
and $\mu_0 = 4\pi \times 10^{-7}\ \si{N.A^{-2}}$ (data sheet). Example: $N = 400$,
$I = \qty{2.0}{A}$, $L = \qty{0.20}{m}$ gives $B = \qty{5.0e-3}{T}$.
\end{refresher}

\begin{formulacard}[combine={$B = \frac{\mu_0 NI}{L}$, $\Phi = BA\cos\theta$ and
    $\varepsilon = -N\frac{\Delta\Phi}{\Delta t}$}]{Emf in a coil wound on a solenoid}
  \cardformula{\lvert\varepsilon_\mathrm{s}\rvert = N_\mathrm{s}\,
    \frac{\mu_0 N_\mathrm{p} A}{L}\,\frac{\lvert\Delta I_\mathrm{p}\rvert}{\Delta t}}
  \begin{symbols}
    \sym{$\varepsilon_\mathrm{s}$}{average emf in the secondary coil}{\si{V}}
    \sym{$N_\mathrm{s}$}{number of turns of the secondary}{no unit}
    \sym{$N_\mathrm{p}$}{number of turns of the primary solenoid}{no unit}
    \sym{$\mu_0$}{magnetic permeability constant,
      $4\pi \times 10^{-7}$ (data sheet)}{\si{N.A^{-2}}}
    \sym{$A$}{cross-section area of the primary}{\si{m^2}}
    \sym{$L$}{length of the primary}{\si{m}}
    \sym{$\Delta I_\mathrm{p}$}{change in the primary current}{\si{A}}
    \sym{$\Delta t$}{time taken for the change}{\si{s}}
  \end{symbols}
  \inwords{the faster the primary current changes, the bigger the emf in the
    secondary.}
  \whentouse{a secondary wound tightly round the middle of a long air-core
    solenoid, so all the primary's flux passes through it.}
  \whennottouse{with an iron core, or with coils far apart (the flux is
    different). A steady primary current gives zero emf.}
  \rearranged{$\Delta t = \dfrac{N_\mathrm{s}\mu_0 N_\mathrm{p} A
    \lvert\Delta I_\mathrm{p}\rvert}{L\lvert\varepsilon_\mathrm{s}\rvert}$ \qquad
    $\lvert\Delta I_\mathrm{p}\rvert = \dfrac{\lvert\varepsilon_\mathrm{s}\rvert L\,
    \Delta t}{N_\mathrm{s}\mu_0 N_\mathrm{p} A}$}
  \tensecond{$N_\mathrm{s} = 50$, $N_\mathrm{p} = 1000$, $A = \qty{1.0e-3}{m^2}$,
    $L = \qty{0.50}{m}$; the current rises by \qty{2.0}{A} in \qty{0.020}{s}: \
    $\lvert\varepsilon_\mathrm{s}\rvert = \dfrac{50 \times 4\pi\times10^{-7}
    \times 1000 \times \num{1.0e-3} \times 2.0}{0.50 \times 0.020}
    = \qty{0.013}{V}$.}
  \derive{only the primary current changes, and $\theta = \ang{0}$.
    \begin{align*}
      B &= \frac{\mu_0 N_\mathrm{p} I_\mathrm{p}}{L} && \text{field inside the primary}\\
      \Delta\Phi &= \Delta B \times A = \frac{\mu_0 N_\mathrm{p} A\,\Delta I_\mathrm{p}}{L}
        && \text{$\Phi = BA$, only $I_\mathrm{p}$ changes}\\
      \lvert\varepsilon_\mathrm{s}\rvert &= N_\mathrm{s}\frac{\lvert\Delta\Phi\rvert}{\Delta t}
        && \text{Faraday, with the secondary's turns}
    \end{align*}}
  \trap{there are two different $N$s. $N_\mathrm{p}$ goes into $B$;
    $N_\mathrm{s}$ goes into Faraday's law. Use the primary's area, not the
    secondary's.}
\end{formulacard}

\begin{workedexample}{A coil wound on a solenoid}\label{ch3:we:twocoils}
A primary solenoid (air core) has 600 turns, length \qty{30}{cm} and radius
\qty{2.0}{cm}. A 40-turn secondary is wound tightly around its middle. The
primary current rises steadily from 0 to \qty{2.5}{A} in \qty{0.010}{s}.
(a) Find the average emf in the secondary. (b) What is the emf once the current
stays at \qty{2.5}{A}? (c) The primary's field points right. Which way does the
secondary's induced field point while the current rises?
\given{$N_\mathrm{p} = 600$, \ $L = \qty{0.30}{m}$, \ $r = \qty{0.020}{m}$, \
  $N_\mathrm{s} = 40$, \ $I$: 0 $\to$ \qty{2.5}{A} in \qty{0.010}{s}. All the
  primary's flux passes through the secondary.}
\find{$\varepsilon$ in the secondary; $\varepsilon$ for a steady current; the
  induced field direction}
\formula{$B = \dfrac{\mu_0 N_\mathrm{p} I}{L}$, \ $\Phi = BA$, \
  $\lvert\varepsilon\rvert = N_\mathrm{s}\dfrac{\lvert\Delta\Phi\rvert}{\Delta t}$}
\substitute{\eqlines{B_\mathrm{f} &= \frac{4\pi\times10^{-7}\,\si{N.A^{-2}} \times 600 \times
    \qty{2.5}{A}}{\qty{0.30}{m}} = \qty{6.283e-3}{T} \quad \text{(final field)}\\
  A &= \pi \times (\qty{0.020}{m})^2 = \qty{1.257e-3}{m^2} \quad \text{(one turn)}\\
  \Delta\Phi &= \qty{6.283e-3}{T} \times \qty{1.257e-3}{m^2} - 0 = \qty{7.90e-6}{Wb}\\
  \lvert\varepsilon\rvert &= \frac{40 \times \qty{7.90e-6}{Wb}}{\qty{0.010}{s}}
    = \qty{3.16e-2}{V} \quad \text{(Faraday)}}}
\answer{(a) \final{$\qty{0.032}{V}$} (\qty{32}{mV}) \ (b) \textbf{zero}: steady
  current, steady flux \ (c) to the \textbf{left}, opposing the growing flux}
\checkstep{a small emf, as expected with an air core. An iron core would give a
  much bigger flux and emf.}
\end{workedexample}

\begin{workedexample}{A magnet dropped through a coil}\label{ch3:we:drop}
A small magnet is dropped from rest, N pole first, through a horizontal coil.
The coil is joined to a data logger (a device that records the emf many times
each second and draws the graph). Describe the emf--time graph and explain its
shape.
\given{the magnet speeds up as it falls; the coil is fixed}
\find{the shape of the $\varepsilon$--$t$ graph, with reasons}
\formula{$\varepsilon = -N\dfrac{\Delta\Phi}{\Delta t}$ and Lenz's law}
\substitute{1. Above the coil the flux increases: a first pulse.\\
  2. At the centre the flux is largest and not changing for an instant:
  $\varepsilon = 0$.\\
  3. Below the coil the flux decreases: a second pulse of opposite sign.\\
  4. The magnet is faster as it leaves, so the flux changes faster: the second
  pulse is taller and narrower.}
\answer{two pulses of opposite sign, the second taller and narrower
  (\cref{ch3:fig:falling}). The area under each pulse (roughly its height
  $\times$ width) is about the same, because each area equals
  $N\lvert\Delta\Phi\rvert$.}
\checkstep{a data logger lets almost no current flow, so the braking is
  negligible. With a closed circuit, the coil repels the magnet above it and
  attracts it below (Lenz). Some of the magnet's gravitational potential energy
  becomes electrical energy and then heat. So it gains less kinetic energy, and
  falls a little slower, than in free fall.}
\begin{figure}
  \centering
  \includegraphics{ch3_falling.pdf}
  \caption{Model: a magnet dropped from rest \qty{6}{cm} above a coil of radius
    \qty{2.5}{cm}.}\label{ch3:fig:falling}
\end{figure}
\end{workedexample}

\begin{workedexample}{Current through a drawn coil}\label{ch3:we:winding}
A vertical coil is joined to an ammeter by leads X and Y
(\cref{ch3:fig:winding}). A bar magnet hangs above the coil, N pole down. It is
lifted straight up. Which way does the current flow through the ammeter?
\begin{figure}
  \centering
  \begin{tikzpicture}[x=1cm, y=1cm]
    % back halves of the turns (faint), then front halves (solid)
    \foreach \i in {0,...,3}
      \draw[line width=0.7pt, hscmuted!70] (0.9,2.375-0.45*\i) -- (-0.9,2.15-0.45*\i);
    \foreach \i in {0,...,4}
      \draw[line width=1.6pt, hscink] (-0.9,2.6-0.45*\i) -- (0.9,2.375-0.45*\i);
    % leads and ammeter
    \draw[line width=0.9pt, hscink] (-0.9,2.6) -- (-0.9,2.95) -- (-2.6,2.95) -- (-2.6,1.95);
    \draw[line width=0.9pt, hscink] (0.9,0.575) -- (0.9,0.15) -- (-2.6,0.15) -- (-2.6,1.25);
    \draw[line width=0.9pt, hscink, fill=white] (-2.6,1.6) circle (0.35);
    \node[font=\small] at (-2.6,1.6) {A};
    \node[font=\small\bfseries, above] at (-1.75,2.95) {X};
    \node[font=\small\bfseries, below] at (-0.85,0.15) {Y};
    % magnet, N pole down, moving up
    \fill[hscforce!22] (-0.3,3.35) rectangle (0.3,3.95);
    \fill[hscvel!22] (-0.3,3.95) rectangle (0.3,4.55);
    \draw[line width=0.8pt, hscink] (-0.3,3.35) rectangle (0.3,4.55);
    \node[font=\small\bfseries] at (0,3.65) {N};
    \node[font=\small\bfseries] at (0,4.25) {S};
    \draw[velocity] (0.65,3.6) -- (0.65,4.45) node[right, text=hscvel] {$v$};
    % key
    \node[font=\footnotesize, anchor=west, align=left] at (1.35,2.2)
      {solid: front of each turn\\ (nearest you)};
    \node[font=\footnotesize, anchor=west, text=hscmuted] at (1.35,1.2) {faint: back};
  \end{tikzpicture}
  \caption{Side view of the coil. Each turn is drawn as a solid front half and a
    faint back half.}\label{ch3:fig:winding}
\end{figure}
\given{the N pole is above the coil and moves away. The solid lines are the front
  of each turn.}
\find{the direction of the current through the ammeter}
\formula{Lenz's law, the right-hand grip rule, then follow the drawn wire}
\substitute{Step 1: field lines leave an N pole, so the field in the coil points
  down. Lifting the magnet makes this downward flux decrease.\\
  Step 2: by Lenz's law the induced field also points down, to oppose the
  decrease. So the top face of the coil becomes an S pole.\\
  Step 3: grip rule: a downward field needs a clockwise current, seen from above.\\
  Step 4: clockwise from above means right to left along the front of the coil.
  Right to left along the solid lines is also upwards. So the current flows up
  the coil, from Y to X.\\
  Step 5: the current leaves the coil at X, so it goes through the ammeter from
  the X side to the Y side.}
\answer{\final{through the ammeter from X to Y}}
\checkstep{the top face is S, so it attracts the rising N pole and pulls against
  the lift: the coil opposes the motion \cmark. Pushing the magnet back down would
  reverse the current.}
\end{workedexample}

\begin{mistakes}
\begin{itemize}
  \item \xmark\ One pulse only, or two pulses of the same sign. \cmark\ In and out
        give opposite signs, with zero at the centre.
  \item \xmark\ ``A steady primary current induces a steady secondary current.''
        \cmark\ Only a \emph{changing} current induces an emf.
  \item \xmark\ ``Copper is not magnetic, so a magnet cannot affect it.''
        \cmark\ Copper conducts; the magnet pushes on the currents induced in it.
  \item \xmark\ Guessing the meter direction from ``clockwise''. \cmark\ Find the
        front of each turn, then follow the current to the leads.
\end{itemize}
\end{mistakes}

\begin{examtip}
``Explain why the galvanometer moves only when the switch is opened or closed''
(3 marks): (1) a changing current gives a changing flux through the secondary;
(2) Faraday: an emf only while the flux changes; (3) Lenz: on and off give
opposite directions. Solid versus slotted rings or tubes: the slot does not stop
the emf. It breaks the current path, so there is little current and little
opposing force. Sketching an emf--time graph (3 marks): two pulses, opposite
signs, zero in between, and both axes labelled with quantity and unit (emf in
\si{V}, time in \si{s}).
\end{examtip}

\begin{quickcheck}
  \item Two coils share an iron core. The primary's switch is closed, kept closed,
        then opened. Describe what the galvanometer in the secondary shows.
  \item A small coil carries a steady current. It is pushed into a larger coil
        joined to a galvanometer, held still, then pulled out. Describe and
        explain the readings.
  \item Explain why an AC supply in the primary gives a continuous emf in the
        secondary, but a DC supply does not.
  \item A solid aluminium plate and a slotted one swing between the poles of a
        magnet. Which stops first? Explain.
  \item A primary solenoid (air core) has 800 turns, length \qty{0.40}{m} and
        cross-section area \qty{5.0}{cm^2}. Its current falls steadily from
        \qty{3.0}{A} to zero in \qty{0.025}{s}. Find the average emf in a 25-turn
        coil wound tightly around its middle.
\end{quickcheck}

%% ===================================================================
\section{Key experiments}\label{ch3:sec:experiments}
\covers{M6.IQ3.02}

These experiments give the evidence the syllabus asks for: magnets, straight
wires, solenoids, metal plates, and one coil near another.

\begin{experiment}{Faraday's induction experiments (1831)}
\aim{To show that a changing flux induces an emf, using a magnet and a coil, a
  moving straight wire, and one coil near another.}
\apparatus{coils of 300--600 turns, bar magnets, a centre-zero galvanometer (or a
  data logger); a U-shaped magnet and a straight wire; a soft iron ring with two
  coils, a battery and a switch; a small coil that fits inside a larger coil
  (\cref{ch3:fig:exp1}).
  \begin{figure}
  \centering
  \begin{tikzpicture}[x=1cm, y=1cm]
    % (a) magnet and coil
    \fill[hscvel!22] (0,0) rectangle (0.6,0.45);
    \fill[hscforce!22] (0.6,0) rectangle (1.2,0.45);
    \draw[line width=0.8pt, hscink] (0,0) rectangle (1.2,0.45);
    \node[font=\footnotesize\bfseries] at (0.3,0.225) {S};
    \node[font=\footnotesize\bfseries] at (0.9,0.225) {N};
    \draw[<->, line width=1.1pt, hscvel] (0.15,0.75) -- (1.05,0.75);
    \foreach \x in {1.6, 1.82, ..., 2.5}
      \draw[line width=1pt, hscink] (\x,0.225) ellipse (0.09 and 0.5);
    \draw[line width=0.9pt, hscink] (1.6,-0.28) -- (1.6,-0.95) -- (1.8,-0.95);
    \draw[line width=0.9pt, hscink] (2.48,-0.28) -- (2.48,-0.95) -- (2.36,-0.95);
    \draw[line width=0.9pt, hscink, fill=white] (2.08,-0.95) circle (0.28);
    \node[font=\footnotesize] at (2.08,-0.95) {G};
    \node[font=\footnotesize] at (1.3,-1.6) {(a) magnet and coil};
    % (b) wire across the pole face
    \begin{scope}[xshift=3.6cm]
      \fill[hscforce!8] (0,-0.8) rectangle (1.9,1.1);
      \draw[line width=0.6pt, hscmuted] (0,-0.8) rectangle (1.9,1.1);
      \foreach \x in {0.25, 0.6, 1.3, 1.65} \foreach \y in {-0.55, -0.05, 0.45, 0.9}
        \node[text=hscfield, font=\footnotesize] at (\x,\y) {$\times$};
      \draw[line width=2.4pt, hscmuted] (0.95,-1.0) -- (0.95,1.3);
      \draw[<->, line width=1.1pt, hscvel] (0.4,0.2) -- (1.5,0.2);
      \draw[line width=0.9pt, hscink] (0.95,1.3) -- (2.6,1.3) -- (2.6,0.48);
      \draw[line width=0.9pt, hscink] (0.95,-1.0) -- (2.6,-1.0) -- (2.6,-0.08);
      \draw[line width=0.9pt, hscink, fill=white] (2.6,0.2) circle (0.28);
      \node[font=\footnotesize] at (2.6,0.2) {G};
      \node[font=\footnotesize, align=center, anchor=south] at (0.95,1.36)
        {gap between the poles\\ of a U-shaped magnet};
      \node[font=\footnotesize, align=center] at (1.25,-1.6) {(b) wire moved across\\ the field (into page)};
    \end{scope}
    % (c) Faraday's iron ring
    \begin{scope}[xshift=10.4cm, yshift=0.1cm]
      \draw[line width=6pt, hscmuted!45] (0,0) circle (0.72);
      \foreach \a in {145, 157, ..., 217}
        \draw[line width=0.9pt, hscink] (\a:0.5) -- (\a:0.94);
      \foreach \a in {-36, -24, ..., 36}
        \draw[line width=0.9pt, hscink] (\a:0.5) -- (\a:0.94);
      \node[font=\footnotesize] at (0,1.05) {iron ring};
      \draw[line width=0.9pt, hscink] (145:0.94) -- (-1.75,0.54) -- (-1.75,0.2);
      \draw[line width=0.9pt, hscink] (215:0.94) -- (-1.75,-0.54) -- (-1.75,-0.25);
      \draw[line width=1.2pt, hscink] (-2.0,0.2) -- (-1.5,0.2);
      \draw[line width=2.2pt, hscink] (-1.88,0.05) -- (-1.62,0.05);
      \fill[hscink] (-1.75,-0.25) circle (1.3pt);
      \draw[line width=0.9pt, hscink] (-1.75,-0.25) -- ++(-60:0.32);
      \draw[line width=0.9pt, hscink] (-1.75,0.05) -- (-1.75,-0.05);
      \node[font=\footnotesize, anchor=east] at (-2.0,0.12) {battery};
      \node[font=\footnotesize, anchor=east] at (-1.85,-0.42) {switch};
      \draw[line width=0.9pt, hscink] (35:0.94) -- (1.7,0.54) -- (1.7,0.28);
      \draw[line width=0.9pt, hscink] (-35:0.94) -- (1.7,-0.54) -- (1.7,-0.28);
      \draw[line width=0.9pt, hscink, fill=white] (1.7,0) circle (0.28);
      \node[font=\footnotesize] at (1.7,0) {G};
      \node[font=\footnotesize, align=center] at (0,-1.7) {(c) Faraday's iron ring:\\ primary (left), secondary (right)};
    \end{scope}
  \end{tikzpicture}
  \caption{Faraday's three induction tests.}\label{ch3:fig:exp1}
  \end{figure}}
\method{\begin{enumerate}
  \item With nothing moving, check that the galvanometer reads zero.
  \item (a) Push the N pole into the coil. Hold it still. Then pull it out.
        Repeat with the S pole. Then repeat faster, with a stronger magnet, and
        with more turns. Finally, move the coil instead of the magnet.
  \item (b) Move the wire across the gap of the U-shaped magnet. Then move it
        along the field lines (from pole to pole).
  \item (c) Close the switch, keep it closed for about \qty{10}{s}, then open it.
  \item (d) Connect the small coil to a battery, so it has a steady current. Push
        it into the larger coil, which is joined to the galvanometer. Hold it
        still, then pull it out.
\end{enumerate}}
\results{Typical results:
  \begin{center}\small
  \begin{tabular}{@{}L{74mm}L{52mm}@{}}
    \toprule
    \colhead{Action} & \colhead{Galvanometer}\\
    \midrule
    magnet and coil still & zero\\
    N pole pushed in & flick left\\
    N pole pulled out & flick right\\
    S pole pushed in & flick right\\
    faster, stronger magnet, or more turns & bigger flick\\
    coil moved instead of the magnet & same results\\
    wire moved across the field & small flick\\
    wire moved along the field & zero\\
    switch closed & flick\\
    switch kept closed & zero\\
    switch opened & flick the other way\\
    small coil (steady current) pushed in & flick\\
    small coil held still inside & zero\\
    small coil pulled out & flick the other way\\
    \bottomrule
  \end{tabular}
  \end{center}}
\conclusion{an emf is induced only while the flux through the circuit changes. It
  grows with the rate of change and the number of turns
  ($\varepsilon = -N\dfrac{\Delta\Phi}{\Delta t}$), and it reverses when the change
  reverses (Lenz). It does not matter whether the magnet or the coil moves; only
  the relative motion matters, and only if it changes the flux. A changing
  current in a nearby coil, or a moving coil with a current, works just like a
  moving magnet.}
\significance{with Ørsted's discovery (1820), this showed that electricity and
  magnetism are linked. Generators, microphones and transformers (\cref{ch4:chap}) all
  use it.}
\validity{\begin{itemize}
  \item \textbf{Validity:} change one variable at a time. Hand speed is hard to
        control, so drop the magnet through the coil from a measured height.
  \item \textbf{Reliability:} repeat each action at least three times; the flicks
        should agree.
  \item \textbf{Accuracy:} flicks are brief, so record the emf with a data logger.
        Before testing Lenz's law, find which way the needle moves for a known
        current (use a cell and a resistor).
\end{itemize}}
\end{experiment}

\begin{experiment}{Eddy currents: falling magnets and swinging plates}
\aim{To show that currents induced in solid metal oppose the motion that causes
  them (evidence for Lenz's law).}
\apparatus{a strong (neodymium) magnet of mass \qty{20}{g}; a non-magnetic metal
  cylinder of the same size and mass; \qty{0.80}{m} tubes of PVC, copper, and
  slotted copper; light gates (sensors that record when an object passes them)
  or a stopwatch; solid and slotted aluminium plates that swing as pendulums; a
  U-shaped magnet (\cref{ch3:fig:tubes}).
  \begin{figure}
  \centering
  \begin{tikzpicture}[x=1cm, y=1cm]
    \foreach \xc/\fillc/\lab in {0/hscmuted!12/PVC, 2.2/brown!35/copper, 4.4/brown!35/slotted copper} {
      \fill[\fillc] (\xc-0.3,0) rectangle (\xc+0.3,2.5);
      \draw[line width=0.8pt, hscink] (\xc-0.3,0) rectangle (\xc+0.3,2.5);
      \fill[hscink] (\xc-0.17,2.1) rectangle (\xc+0.17,2.35);
      \draw[velocity] (\xc+0.55,2.25) -- (\xc+0.55,1.5);
      \node[font=\footnotesize, align=center, anchor=north] at (\xc,-0.1) {\lab};}
    \draw[line width=1.6pt, white] (4.4,0.05) -- (4.4,2.45);
    \foreach \y in {1.9, 0.3}
      \draw[line width=0.6pt, hscink, fill=hscbrand!15] (-0.75,\y-0.1) rectangle (-0.45,\y+0.1);
    \node[font=\footnotesize, anchor=east] at (-0.8,1.9) {light gate};
    \node[font=\footnotesize, anchor=east] at (-0.8,0.3) {light gate};
    \draw[<->, line width=0.6pt, hscmuted] (5.3,0) -- (5.3,2.5)
      node[midway, right, text=hscink, font=\footnotesize] {\qty{0.80}{m}};
    \draw[guide] (-0.17,2.25) -- (-0.6,2.6) node[left, font=\footnotesize] {magnet};
    % swinging plates between the poles of a U-shaped magnet (B into the page)
    \foreach \xc/\lab in {7.2/{solid:\\ stops quickly}, 9.4/{slotted:\\ swings longer}} {
      \fill[hscforce!10] (\xc-0.75,0.15) rectangle (\xc+0.75,1.25);
      \draw[line width=0.6pt, hscmuted] (\xc-0.75,0.15) rectangle (\xc+0.75,1.25);
      \foreach \x in {-0.57, 0.57} \foreach \y in {0.35, 0.7, 1.05}
        \node[text=hscfield, font=\footnotesize] at (\xc+\x,\y) {$\times$};
      \fill[hscink] (\xc,2.5) circle (1.6pt);
      \draw[line width=1.0pt, hscink] (\xc,2.5) -- (\xc,1.15);
      \draw[<->, line width=0.9pt, hscvel] ($(\xc,2.5)+(240:0.85)$) arc (240:300:0.85);
      \node[font=\footnotesize, align=center, anchor=north] at (\xc,-0.1) {\lab};}
    % solid plate
    \filldraw[fill=brown!35, draw=hscink, line width=0.8pt] (6.8,0.35) rectangle (7.6,1.15);
    % slotted (comb) plate: three slots cut up from the bottom edge
    \filldraw[fill=brown!35, draw=hscink, line width=0.8pt]
      (9.0,1.15) -- (9.8,1.15) -- (9.8,0.35) -- (9.68,0.35) -- (9.68,1.0)
      -- (9.56,1.0) -- (9.56,0.35) -- (9.46,0.35) -- (9.46,1.0) -- (9.34,1.0)
      -- (9.34,0.35) -- (9.24,0.35) -- (9.24,1.0) -- (9.12,1.0) -- (9.12,0.35)
      -- (9.0,0.35) -- cycle;
  \end{tikzpicture}
  \caption{Left: a magnet dropped through three tubes. Right: plates swinging
    between magnet poles (shaded).}\label{ch3:fig:tubes}
  \end{figure}}
\method{\begin{enumerate}
  \item Release the magnet from rest at the top of the PVC tube; record the fall
        time. Repeat with the copper and slotted copper tubes.
  \item Drop the non-magnetic cylinder through the copper tube.
  \item Do three trials of each and find the mean.
  \item Swing a solid aluminium plate, then a slotted one, between the poles of a
        strong magnet. Count the swings before each stops.
\end{enumerate}}
\results{Sample class results (fall times through the tube):
  \begin{center}\small
  \begin{tabular}{@{}L{60mm}C{14mm}C{14mm}C{14mm}C{15mm}@{}}
    \toprule
    \colhead{Object and tube} & \colhead{Trial 1} & \colhead{Trial 2}
      & \colhead{Trial 3} & \colhead{Mean}\\
    \midrule
    magnet, PVC tube & \qty{0.41}{s} & \qty{0.40}{s} & \qty{0.42}{s} & \qty{0.41}{s}\\
    non-magnetic cylinder, copper tube & \qty{0.40}{s} & \qty{0.42}{s} & \qty{0.41}{s} & \qty{0.41}{s}\\
    magnet, slotted copper tube & \qty{0.55}{s} & \qty{0.58}{s} & \qty{0.54}{s} & \qty{0.56}{s}\\
    magnet, copper tube & \qty{3.52}{s} & \qty{3.61}{s} & \qty{3.58}{s} & \qty{3.57}{s}\\
    \bottomrule
  \end{tabular}
  \end{center}
  Free fall through \qty{0.80}{m}: from $s = ut + \frac12 at^2$ with $u = 0$,
  $t = \sqrt{2s/g} = \sqrt{2 \times \qty{0.80}{m} \div \qty{9.8}{m.s^{-2}}}
  = \qty{0.40}{s}$.
  The solid plate stops within a few swings; the slotted plate swings many more
  times.}
\conclusion{only a magnet in a conducting tube is slowed. The falling magnet
  changes the flux through each ring of copper, inducing eddy currents (Faraday).
  By Lenz's law the copper below repels it and the copper above attracts it: an
  upward force. The magnet soon reaches a steady speed of about
  $\qty{0.80}{m} \div \qty{3.57}{s} = \qty{0.224}{m.s^{-1}}$. The \qty{20}{g}
  (\qty{0.020}{kg}) magnet loses
  $\Delta U = mg\Delta h = \qty{0.020}{kg} \times \qty{9.8}{m.s^{-2}} \times
  \qty{0.80}{m} = \qty{0.157}{J}$ of gravitational potential energy. It gains only
  $K = \frac12 mv^2 = \frac12 \times \qty{0.020}{kg} \times
  (\qty{0.224}{m.s^{-1}})^2 = \qty{5.0e-4}{J}$ of kinetic energy. Fraction turned
  into heat $= (0.157 - 0.0005) \div 0.157 = 0.997$, so over \qty{99}{\percent}
  becomes heat in the tube. The slot breaks the current loops, so the braking is
  much weaker. The same happens with the plates.}
\significance{clear evidence for Lenz's law and energy conservation; the idea
  behind magnetic brakes (\cref{ch5:chap}) and layered transformer cores (\cref{ch4:chap}).
  HSC questions sometimes stand the copper tube on a balance. The tube pushes up
  on the falling magnet. By Newton's third law, the magnet pushes down on the
  tube with an equal force. At the steady speed, this push equals the magnet's
  weight. So the balance reads the same as if the magnet were resting on the
  tube.}
\validity{\begin{itemize}
  \item \textbf{Validity:} same magnet, drop height and tube length. The PVC tube
        and the non-magnetic cylinder are controls: slowing needs both a magnet and
        a conductor.
  \item \textbf{Reliability:} three or more trials; these agree within \qty{0.1}{s}.
  \item \textbf{Accuracy:} stopwatch reaction time (about \qty{0.2}{s}) is half the
        free-fall time, so use light gates or slow-motion video.
\end{itemize}}
\end{experiment}

%% ===================================================================
\begin{summarybox}{\Cref{ch3:chap} in five lines}
\begin{itemize}
  \item $\Phi = BA\cos\theta$, with $\theta$ from the normal. It changes when $B$,
        $A$ or $\theta$ changes.
  \item Faraday: $\varepsilon = -N\dfrac{\Delta\Phi}{\Delta t}$. Only a changing flux
        induces an emf; a current also needs a closed circuit.
  \item Lenz: the induced current opposes the change in flux. Work is done against
        it, so energy is conserved.
  \item Moving rod: $\varepsilon = lvB$ (not on the sheet); the force on the
        induced current opposes the motion.
  \item Magnets in solenoids, eddy currents in plates and two nearby coils all
        follow the same two laws.
\end{itemize}
\end{summarybox}
